\chapter{源代码清单}

本系统分为后端（Python/Flask）和前端（Flutter/Dart）两大部分。以下选取核心模块的源代码进行展示。

\section{后端源代码}

\subsection{主程序入口 (back/main.py)}
\begin{lstlisting}[style=codestyle, language=Python, caption={后端主程序入口}]
from flask import Flask, jsonify, request
from flask_cors import CORS
from config import config
from services.firebase_service import FirebaseService
from middleware.logging import setup_logging
from utils.exceptions import register_error_handlers
from scheduler import get_scheduler
import os

from api.auth import auth_bp
from api.data import data_bp
from api.analysis import analysis_bp
from api.history import history_bp
from api.optimization import optimization_bp

def create_app(config_name=None):
    """应用工厂函数"""
    app = Flask(__name__)
    
    if config_name is None:
        config_name = os.getenv('FLASK_ENV', 'development')
    app.config.from_object(config[config_name])
    config[config_name].init_app(app)
    
    CORS(app, 
         origins=app.config['CORS_ORIGINS'],
         supports_credentials=True)
    
    setup_logging(app)
    register_error_handlers(app)
    
    try:
        FirebaseService.initialize(project_id=app.config['GCP_PROJECT_ID'])
    except Exception as e:
        print(f"Firebase Init Failed: {e}")
    
    app.register_blueprint(auth_bp)
    app.register_blueprint(data_bp)
    app.register_blueprint(analysis_bp)
    app.register_blueprint(history_bp)
    app.register_blueprint(optimization_bp)
    
    # GAE/Cloud Run Cron Handler
    @app.route('/tasks/fetch-data', methods=['GET'])
    def trigger_fetch_data():
        if os.getenv('FLASK_ENV') != 'development' and request.headers.get('X-Appengine-Cron') != 'true':
            return jsonify({'error': 'Unauthorized'}), 403
        try:
            get_scheduler().fetch_data_job()
            return jsonify({'status': 'success'}), 200
        except Exception as e:
            return jsonify({'error': str(e)}), 500

    return app

app = create_app()

if __name__ == '__main__':
    port = int(os.getenv('PORT', 8080))
    app.run(host='0.0.0.0', port=port, debug=True)
\end{lstlisting}

\subsection{系统配置 (back/config.py)}
\begin{lstlisting}[style=codestyle, language=Python, caption={后端配置定义}]
import os
from dotenv import load_dotenv

load_dotenv()

class Config:
    DEBUG = False
    GCP_PROJECT_ID = os.getenv('GCP_PROJECT_ID', 'data-science-44398')
    PRICE_SCHEDULE = {
        'valley': 0.3,
        'normal': 0.6,
        'peak': 1.0,
        'valley_hours_list': [0, 1, 2, 3, 4, 5, 6, 7, 22, 23],
        'normal_hours_list': [8, 9, 10, 11, 12, 13, 14, 15, 16, 17],
        'peak_hours_list': [18, 19, 20, 21]
    }
    CORS_ORIGINS = ['https://data-science-44398.web.app']

class DevelopmentConfig(Config):
    DEBUG = True
    CORS_ORIGINS = Config.CORS_ORIGINS + ['http://localhost:*']

class ProductionConfig(Config):
    DEBUG = False

config = {
    'development': DevelopmentConfig,
    'production': ProductionConfig,
    'default': DevelopmentConfig
}
\end{lstlisting}

\subsection{能源预测服务 (back/services/ml\_service.py)}
\begin{lstlisting}[style=codestyle, language=Python, caption={机器学习预测服务}]
"""
机器学习服务模块 - 能源负载预测
ML Service Module for Energy Load Prediction

支持多种模型：RandomForest、LightGBM、XGBoost
自动模型选择和超参数优化
"""

import pandas as pd
import numpy as np
import joblib
import os
import tempfile
from pathlib import Path
from typing import List, Dict, Optional, Union, Any
from datetime import datetime, timedelta
from sklearn.ensemble import RandomForestRegressor, VotingRegressor
from sklearn.model_selection import train_test_split, TimeSeriesSplit, RandomizedSearchCV
from sklearn.metrics import mean_absolute_error, mean_squared_error, r2_score
import warnings

warnings.filterwarnings('ignore')

# 可选依赖：LightGBM 和 XGBoost
try:
    from lightgbm import LGBMRegressor
    LIGHTGBM_AVAILABLE = True
except ImportError:
    LIGHTGBM_AVAILABLE = False
    print(" LightGBM 未安装，将使用 RandomForest")

try:
    from xgboost import XGBRegressor
    XGBOOST_AVAILABLE = True
except ImportError:
    XGBOOST_AVAILABLE = False
    print(" XGBoost 未安装，将使用 RandomForest")


from config import Config

class EnergyPredictor:
    """
    能源负载预测器
    
    使用随机森林模型预测未来24小时的能源负载
    支持模型训练、保存、加载和推理
    """
    
    def __init__(self, model_path: str = None):
        """
        初始化预测器
        
        Args:
            model_path: 本地兜底模型路径（可选），主要从 Firebase Storage 加载
        """
        # 获取项目根目录
        self.script_dir = Path(__file__).parent  # services 目录
        self.back_dir = self.script_dir.parent  # back 目录
        
        # Firebase Storage 模型路径（主要存储位置）
        self.firebase_model_path = 'models/rf_model.joblib'
        
        # 本地兜底模型路径（部署包中自带的模型，仅用于首次加载）
        if model_path:
            self.local_model_path = Path(model_path)
        else:
            self.local_model_path = self.back_dir / 'models' / 'rf_model.joblib'
        
        # 初始化 StorageService
        from services.storage_service import StorageService
        self.storage_service = StorageService()
        
        # 初始化模型
        self.model: Optional[RandomForestRegressor] = None
        
        # 基础特征列表（向后兼容）
        self.base_feature_columns = [
            'Hour', 'DayOfWeek', 'Temperature', 'Price',
            'Lag_1h', 'Lag_24h', 'Lag_168h',
            'Rolling_Mean_6h', 'Rolling_Std_6h', 'Rolling_Mean_24h',
            'Temp_x_Hour', 'Lag24_x_DayOfWeek'
        ]
        
        # 增强特征列表（新增的时间特征）
        self.enhanced_feature_columns = [
            # 基础时间特征
            'Month', 'Season', 'IsWeekend', 'IsHoliday', 'DayOfMonth', 'WeekOfYear',
            # 增强交互特征
            'Temp_x_Season', 'Lag24_x_IsWeekend', 'Hour_x_IsHoliday',
            # 周期编码特征
            'Month_Sin', 'Month_Cos', 'Hour_Sin', 'Hour_Cos'
        ]
        
        # 实际使用的特征列表（初始化为基础特征）
        self.feature_columns = self.base_feature_columns.copy()
        self.target_column = 'Site_Load'
        
        print(f" Firebase Storage 模型路径: {self.firebase_model_path}")
        print(f" 本地兜底模型路径: {self.local_model_path}")
    
    def _load_feature_columns_from_metadata(self):
        """
        从模型元数据中加载特征列表
        确保预测时使用与训练时相同的特征
        """
        try:
            metadata = self.get_model_metadata()
            if metadata and 'feature_engineering' in metadata:
                fe_info = metadata['feature_engineering']
                if fe_info.get('use_enhanced', False):
                    # 模型使用了增强特征，更新特征列表
                    enhanced_list = fe_info.get('enhanced_features_list', [])
                    if enhanced_list:
                        self.feature_columns = self.base_feature_columns.copy()
                        self.feature_columns.extend(enhanced_list)
                        print(f"    已从元数据恢复特征列表: {len(self.feature_columns)} 个特征")
                        print(f"     (基础: {len(self.base_feature_columns)}, 增强: {len(enhanced_list)})")
                else:
                    # 模型只使用基础特征
                    self.feature_columns = self.base_feature_columns.copy()
                    print(f"    模型使用基础特征: {len(self.feature_columns)} 个")
            else:
                print(f"     未找到特征元数据，使用默认基础特征")
                
            # 加载训练配置 (Log Transform)
            if metadata and 'training_config' in metadata:
                config = metadata['training_config']
                self.use_log_transform = config.get('use_log_transform', False)
                if self.use_log_transform:
                    print(f"    模型使用 Log1p 变换，预测时将自动还原")
            else:
                self.use_log_transform = False
                
        except Exception as e:
            print(f"     加载特征元数据失败: {e}，使用默认基础特征")
            self.use_log_transform = False
    
    def _get_model_type_name(self) -> str:
        """获取当前模型的类型名称"""
        if self.model is None:
            return 'Unknown'
        
        model_class = type(self.model).__name__
        name_map = {
            'RandomForestRegressor': 'Random Forest Regressor',
            'LGBMRegressor': 'LightGBM Regressor',
            'XGBRegressor': 'XGBoost Regressor'
        }
        return name_map.get(model_class, model_class)
    
    def _create_model(self, model_type: str, n_estimators: int = 100, random_state: int = 42):
        """
        创建指定类型的模型
        
        Args:
            model_type: 模型类型 ('randomforest', 'lightgbm', 'xgboost')
            n_estimators: 树的数量
            random_state: 随机种子
            
        Returns:
            (model, hyperparameters) 元组
        """
        model_type = model_type.lower()
        
        if model_type == 'lightgbm' and LIGHTGBM_AVAILABLE:
            model = LGBMRegressor(
                n_estimators=n_estimators,
                learning_rate=0.05,
                max_depth=15,
                num_leaves=31,
                subsample=0.8,
                colsample_bytree=0.8,
                random_state=random_state,
                n_jobs=-1,
                verbose=-1
            )
            params = {
                'n_estimators': n_estimators,
                'learning_rate': 0.05,
                'max_depth': 15,
                'num_leaves': 31
            }
        elif model_type == 'xgboost' and XGBOOST_AVAILABLE:
            model = XGBRegressor(
                n_estimators=n_estimators,
                learning_rate=0.05,
                max_depth=10,
                subsample=0.8,
                colsample_bytree=0.8,
                random_state=random_state,
                n_jobs=-1,
                verbosity=0
            )
            params = {
                'n_estimators': n_estimators,
                'learning_rate': 0.05,
                'max_depth': 10
            }
        else:
            # 默认使用 RandomForest
            model = RandomForestRegressor(
                n_estimators=n_estimators,
                max_depth=20,
                min_samples_split=5,
                min_samples_leaf=2,
                random_state=random_state,
                n_jobs=-1,
                verbose=0
            )
            params = {
                'n_estimators': n_estimators,
                'max_depth': 20,
                'min_samples_split': 5,
                'min_samples_leaf': 2
            }
        
        return model, params
    
    def _tune_model(self, model_type: str, X_train, y_train, cv_split, n_iter: int = 20, random_state: int = 42):
        """
        使用 RandomizedSearchCV 对指定模型进行超参数调优
        """
        model_type = model_type.lower()
        estimator = None
        param_dist = {}
        
        if model_type == 'lightgbm' and LIGHTGBM_AVAILABLE:
            estimator = LGBMRegressor(random_state=random_state, n_jobs=-1, verbose=-1)
            param_dist = {
                'n_estimators': [100, 200, 300, 500],
                'learning_rate': [0.01, 0.03, 0.05, 0.1],
                'num_leaves': [31, 50, 70, 100],
                'max_depth': [10, 15, 20, -1],
                'subsample': [0.7, 0.8, 0.9],
                'colsample_bytree': [0.7, 0.8, 0.9]
            }
        elif model_type == 'xgboost' and XGBOOST_AVAILABLE:
            estimator = XGBRegressor(random_state=random_state, n_jobs=-1, verbosity=0)
            param_dist = {
                'n_estimators': [100, 200, 300, 500],
                'learning_rate': [0.01, 0.03, 0.05, 0.1],
                'max_depth': [3, 6, 10, 15],
                'subsample': [0.7, 0.8, 0.9],
                'colsample_bytree': [0.7, 0.8, 0.9],
                'gamma': [0, 0.1, 0.2]
            }
        else:
            # RandomForest
            estimator = RandomForestRegressor(random_state=random_state, n_jobs=-1)
            param_dist = {
                'n_estimators': [100, 200, 300],
                'max_depth': [10, 20, 30, None],
                'min_samples_split': [2, 5, 10],
                'min_samples_leaf': [1, 2, 4],
                'max_features': ['sqrt', 'log2', None]
            }
            
        print(f"    正在为 {model_type} 搜索最佳参数 (iter={n_iter})...")
        
        search = RandomizedSearchCV(
            estimator=estimator,
            param_distributions=param_dist,
            n_iter=n_iter,
            cv=cv_split,
            scoring='neg_mean_absolute_error',
            random_state=random_state,
            n_jobs=-1,
            verbose=0
        )
        
        search.fit(X_train, y_train)
        print(f"       最佳参数: {search.best_params_}")
        return search.best_estimator_, search.best_params_

    def _auto_select_best_model(
        self, 
        X_train, y_train, 
        X_test, y_test, 
        random_state: int = 42,
        use_time_series_cv: bool = True,
        n_splits: int = 5,
        perform_tuning: bool = True  # 新增参数
    ) -> tuple:
        """
        自动选择最佳模型（支持时间序列交叉验证）
        
        比较多种模型配置，选择测试集 MAE 最低的
        
        Args:
            X_train, y_train: 训练数据
            X_test, y_test: 测试数据
            random_state: 随机种子
            use_time_series_cv: 是否使用时间序列交叉验证（默认 True）
            n_splits: 交叉验证折数（默认 5）
            
        Returns:
            (best_model, best_params, selection_info) 元组
        """
        candidates = []
        results = {}
        cv_details = {}
        
        if perform_tuning:
            # 如果启用调优，此时我们不使用预设的配置，而是针对每个模型类型运行搜索
            # 但为了保持代码结构，我们构建一个待调优类型列表
            # 这里的 (DisplayName, type, init_n_estimators) 仅用于元组解包，n_estimators 会被忽略
            tuning_targets = [('RandomForest', 'randomforest', 100)]
            if LIGHTGBM_AVAILABLE:
                tuning_targets.append(('LightGBM', 'lightgbm', 100))
            if XGBOOST_AVAILABLE:
                tuning_targets.append(('XGBoost', 'xgboost', 100))
            
            model_configs = tuning_targets # 替换为调优目标
        else:
            # 原有的固定配置模式
            model_configs = [
                ('RandomForest', 'randomforest', 150),
                ('RandomForest_200', 'randomforest', 200),
            ]
            
            # 如果 LightGBM 可用，添加到候选
            if LIGHTGBM_AVAILABLE:
                model_configs.extend([
                    ('LightGBM', 'lightgbm', 200),
                    ('LightGBM_300', 'lightgbm', 300),
                ])
            
            # 如果 XGBoost 可用，添加到候选
            if XGBOOST_AVAILABLE:
                model_configs.extend([
                    ('XGBoost', 'xgboost', 200),
                    ('XGBoost_300', 'xgboost', 300),
                ])
        
        print(f"   评估 {len(model_configs)} 种模型配置...")
        if use_time_series_cv:
            print(f"    使用 TimeSeriesSplit 交叉验证 ({n_splits} 折)")
        
        baseline_mae = None
        best_mae = float('inf')
        best_model = None
        best_params = None
        best_name = None
        
        # 【修复】时间序列交叉验证只在训练集上进行，测试集保持独立用于最终评估
        # 这避免了测试集参与模型选择导致的数据泄漏
        if use_time_series_cv:
            # CV 只在训练集上做，不再把测试集拼进去
            tscv = TimeSeriesSplit(n_splits=n_splits)
            print(f"    TimeSeriesSplit CV 仅在训练集上进行 ({len(X_train)} 样本)")
        
        # 存储每种类型的最佳模型，用于集成
        best_estimators = {}
        
        for name, model_type, n_estimators in model_configs:
            try:
                model = None
                params = None
                
                if perform_tuning:
                    # 使用超参数搜索
                    # 注意：复用 tscv 作为 cv_split
                    # 如果未启用 TimeSeriesCV，则使用 KFold 或者默认的 5折
                    cv_to_use = tscv if use_time_series_cv else 5
                    
                    # 调优后的名字加上 "Tuned" 后缀
                    name = f"{name}_Tuned"
                    model, params = self._tune_model(model_type, X_train, y_train, cv_to_use, n_iter=15, random_state=random_state)
                else:
                    # 使用固定配置
                    print(f"   - 训练 {name}...", end=' ')
                    model, params = self._create_model(model_type, n_estimators, random_state)
                
                if use_time_series_cv:
                    # 【修复】时间序列交叉验证只在训练集上进行
                    cv_scores = []
                    for train_idx, val_idx in tscv.split(X_train):
                        X_cv_train, X_cv_val = X_train.iloc[train_idx], X_train.iloc[val_idx]
                        y_cv_train, y_cv_val = y_train.iloc[train_idx], y_train.iloc[val_idx]
                        
                        model_cv, _ = self._create_model(model_type, n_estimators, random_state)
                        model_cv.fit(X_cv_train, y_cv_train)
                        y_cv_pred = model_cv.predict(X_cv_val)
                        
                        # 如果启用了对数变换，需要还原预测值
                        if hasattr(self, 'use_log_transform') and self.use_log_transform:
                            y_cv_val = np.expm1(y_cv_val)
                            y_cv_pred = np.expm1(y_cv_pred)
                            
                        cv_scores.append(mean_absolute_error(y_cv_val, y_cv_pred))
                    
                    # 计算交叉验证平均分
                    cv_mae = np.mean(cv_scores)
                    cv_std = np.std(cv_scores)
                    
                    # 使用完整训练数据重新训练最终模型
                    model.fit(X_train, y_train)
                    
                    # 在独立的未来测试集上评估（严格未见数据）
                    y_pred = model.predict(X_test)
                    
                    # 如果启用了对数变换，需要还原预测值
                    if hasattr(self, 'use_log_transform') and self.use_log_transform:
                         # y_test 是原始值（如果不改动 y_test 的话），这里需要确认 split 时 y_test 是否被 log
                         # 我们的策略是：只对训练集的 y 进行 log，测试集的 y 保持原样或者反变换
                         # 为统一处理，假设 fit 时 y_train 是 log 后的，predict 输出 log 后的
                         y_pred = np.expm1(y_pred)
                         # y_test 在 split 前如果没变，那就是原始值。
                         # 为了安全，我们稍后在 split 处做处理。
                    
                    test_mae = mean_absolute_error(y_test, y_pred)
                    test_rmse = np.sqrt(mean_squared_error(y_test, y_pred))
                    
                    # 使用交叉验证 MAE 作为选择依据（更可靠）
                    mae = cv_mae
                    rmse = test_rmse
                    
                    cv_details[name] = {
                        'cv_mae_mean': round(cv_mae, 2),
                        'cv_mae_std': round(cv_std, 2),
                        'cv_scores': [round(s, 2) for s in cv_scores],
                        'test_mae': round(test_mae, 2)  # 记录独立测试集的 MAE
                    }
                    
                    print(f"CV_MAE={cv_mae:.2f}±{cv_std:.2f} kW, Test_MAE={test_mae:.2f} kW (独立未来数据)")
                else:
                    # 原有的简单训练-测试拆分
                    model.fit(X_train, y_train)
                    y_pred = model.predict(X_test)
                    
                    if hasattr(self, 'use_log_transform') and self.use_log_transform:
                        y_pred = np.expm1(y_pred)
                    
                    mae = mean_absolute_error(y_test, y_pred)
                    rmse = np.sqrt(mean_squared_error(y_test, y_pred))
                    print(f"MAE={mae:.2f} kW")
                
                results[name] = {'mae': mae, 'rmse': rmse, 'model_type': model_type}
                candidates.append(name)
                
                # 记录基准（第一个 RandomForest）
                if baseline_mae is None:
                    baseline_mae = mae
                
                # 更新最佳
                if mae < best_mae:
                    best_mae = mae
                    best_model = model
                    best_params = params
                    best_name = name
                    
            except Exception as e:
                print(f"失败: {str(e)}")
                continue
            
            # 记录该类型的最佳模型用于集成
            if model_type not in best_estimators or mae < best_estimators[model_type]['mae']:
                best_estimators[model_type] = {
                    'model': model,
                    'mae': mae
                }
        
        # 尝试集成学习 (VotingRegressor)
        if len(best_estimators) >= 2:
            try:
                print(f"    尝试集成学习 (VotingRegressor)...", end=' ')
                # 创建 VotingRegressor
                estimators_list = [(k, v['model']) for k, v in best_estimators.items()]
                voting_model = VotingRegressor(estimators=estimators_list, n_jobs=-1)
                voting_name = "Voting_Ensemble"
                
                # 评估 Voting 模型
                if use_time_series_cv:
                    cv_scores = []
                    for train_idx, val_idx in tscv.split(X_train):
                        X_cv_train, X_cv_val = X_train.iloc[train_idx], X_train.iloc[val_idx]
                        y_cv_train, y_cv_val = y_train.iloc[train_idx], y_train.iloc[val_idx]
                        
                        from sklearn.base import clone
                        model_cv = clone(voting_model)
                        model_cv.fit(X_cv_train, y_cv_train)
                        y_cv_pred = model_cv.predict(X_cv_val)
                        cv_scores.append(mean_absolute_error(y_cv_val, y_cv_pred))
                    
                    cv_mae = np.mean(cv_scores)
                    cv_std = np.std(cv_scores)
                    
                    voting_model.fit(X_train, y_train)
                    y_pred = voting_model.predict(X_test)
                    
                    if hasattr(self, 'use_log_transform') and self.use_log_transform:
                         y_pred = np.expm1(y_pred)
                         
                    test_mae = mean_absolute_error(y_test, y_pred)
                    test_rmse = np.sqrt(mean_squared_error(y_test, y_pred))
                    
                    mae = cv_mae
                    rmse = test_rmse
                    
                    cv_details[voting_name] = {
                        'cv_mae_mean': round(cv_mae, 2),
                        'cv_mae_std': round(cv_std, 2),
                        'cv_scores': [round(s, 2) for s in cv_scores],
                        'test_mae': round(test_mae, 2)
                    }
                    print(f"CV_MAE={cv_mae:.2f}±{cv_std:.2f} kW, Test_MAE={test_mae:.2f} kW")
                else:
                    voting_model.fit(X_train, y_train)
                    y_pred = voting_model.predict(X_test)
                    
                    if hasattr(self, 'use_log_transform') and self.use_log_transform:
                        y_pred = np.expm1(y_pred)
                    
                    mae = mean_absolute_error(y_test, y_pred)
                    rmse = np.sqrt(mean_squared_error(y_test, y_pred))
                    print(f"MAE={mae:.2f} kW")
                
                results[voting_name] = {'mae': mae, 'rmse': rmse, 'model_type': 'voting'}
                candidates.append(voting_name)
                
                if mae < best_mae:
                    best_mae = mae
                    best_model = voting_model
                    best_params = {'estimators': [k for k in best_estimators.keys()]}
                    best_name = voting_name
                    print(f"    集成模型表现最佳!")
                    
            except Exception as e:
                print(f"    集成学习失败: {str(e)}")

        # 计算相对基准的提升
        improvement = 'N/A'
        if baseline_mae and baseline_mae > 0:
            improvement_pct = (baseline_mae - best_mae) / baseline_mae * 100
            improvement = f"{improvement_pct:.1f}%"
        
        print(f"\n    最佳模型: {best_name} (MAE={best_mae:.2f} kW)")
        if improvement != 'N/A' and improvement != '0.0%':
            print(f"    相比基准提升: {improvement}")
        
        selection_info = {
            'winner': best_name,
            'candidates_evaluated': candidates,
            'all_scores': {k: {'mae': round(v['mae'], 2), 'rmse': round(v['rmse'], 2)} 
                          for k, v in results.items()},
            'improvement': improvement,
            'validation_method': 'TimeSeriesSplit' if use_time_series_cv else 'HoldOut',
            'cv_folds': n_splits if use_time_series_cv else None,
            'cv_details': cv_details if use_time_series_cv else None
        }
        
        return best_model, best_params, selection_info
    
    def _get_price(self, hour: int) -> float:
        """
        根据小时返回峰谷电价 (从配置读取)
        
        Args:
            hour: 小时 (0-23)
            
        Returns:
            电价 (元/kWh)
        """
        schedule = Config.PRICE_SCHEDULE
        
        if hour in schedule['peak_hours_list']:
            return schedule['peak']
        elif hour in schedule['normal_hours_list']:
            return schedule['normal']
        else:
            return schedule['valley']
    
    def _save_model_metadata(self, metadata: dict) -> bool:
        """
        保存模型元数据到 Firebase Storage (JSON 文件)
        
        Args:
            metadata: 模型元数据字典
            
        Returns:
            是否保存成功
        """
        try:
            import json
            from datetime import datetime
            
            # 添加更新时间戳
            metadata['updated_at'] = datetime.now().isoformat()
            
            # 转换为 JSON 字符串
            json_data = json.dumps(metadata, indent=2, ensure_ascii=False)
            
            # 上传到 Storage (使用临时文件确保 Content-Type 正确)
            import tempfile
            metadata_path = 'models/model_metadata.json'
            
            with tempfile.NamedTemporaryFile(mode='w', suffix='.json', delete=False, encoding='utf-8') as f:
                f.write(json_data)
                temp_path = f.name
            
            try:
                with open(temp_path, 'rb') as f:
                    self.storage_service.bucket.blob(metadata_path).upload_from_file(
                        f, 
                        content_type='application/json'
                    )
            finally:
                import os
                os.unlink(temp_path)
            
            print(f"    模型元数据已保存到 Firebase Storage: {metadata_path}")
            return True
            
        except Exception as e:
            print(f"    保存模型元数据失败: {str(e)}")
            return False
    
    @staticmethod
    def get_model_metadata() -> Optional[dict]:
        """
        从 Firebase Storage 获取模型元数据 (JSON 文件)
        
        Returns:
            模型元数据字典，如果不存在返回 None
        """
        try:
            import json
            from services.storage_service import StorageService
            
            # 创建 Storage 服务实例
            storage = StorageService()
            
            # 下载元数据 JSON
            metadata_path = 'models/model_metadata.json'
            json_bytes = storage.download_file(metadata_path)
            
            # 解析 JSON
            metadata = json.loads(json_bytes.decode('utf-8'))
            return metadata
                
        except Exception as e:
            print(f"获取模型元数据失败: {str(e)}")
            return None
    
    def _is_gae_environment(self) -> bool:
        """检测是否在 GAE 环境"""
        return bool(os.getenv('GAE_ENV') or os.getenv('K_SERVICE'))

    def train_model(
        self, 
        data_path: str = None,
        n_estimators: int = 100,
        test_size: float = 0.2,
        random_state: int = 42,
        use_firebase_storage: bool = True,
        auto_select_model: bool = True,
        model_type: str = 'auto',
        use_enhanced_features: bool = True,
        use_time_series_cv: bool = True,
        cv_folds: int = 5,
        use_log_transform: bool = True,
        remove_outliers: bool = True,
        tune_hyperparameters: bool = True  # 新增控制开关 (默认开启)
    ) -> Dict[str, float]:
        """
        训练预测模型（支持自动模型选择、增强特征和时间序列交叉验证）
        
        Args:
            data_path: 数据文件路径
            n_estimators: 树的数量
            test_size: 测试集比例
            random_state: 随机种子
            use_firebase_storage: 是否从 Firebase Storage 下载数据
            auto_select_model: 是否自动选择最佳模型
            model_type: 指定模型类型
            use_enhanced_features: 是否使用增强特征
            use_time_series_cv: 是否使用时间序列交叉验证
            cv_folds: 交叉验证折数
            use_log_transform: 是否对目标变量进行 Log1p 变换 (默认 True)
            remove_outliers: 是否自动过滤异常值 (默认 True)
            tune_hyperparameters: 是否启用超参数自动搜索 (默认 True) 
            
        Returns:
            包含评估指标的字典
        """
        print("\n" + "="*80)
        print(" 开始训练能源负载预测模型")
        print("="*80 + "\n")
        
        temp_data_path = None
        
        try:
            # 从 Firebase Storage 下载数据
            if use_firebase_storage:
                print(" 从 Firebase Storage 下载训练数据...")
                from services.storage_service import StorageService
                
                storage_service = StorageService()
                firebase_path = data_path or 'data/processed/cleaned_energy_data_all.csv'
                
                temp_data_path = storage_service.download_to_temp(firebase_path)
                
                if temp_data_path is None:
                    raise FileNotFoundError(f"无法从 Firebase Storage 下载数据: {firebase_path}")
                
                data_path = temp_data_path
                print(f"    数据已下载到: {data_path}")
            else:
                # 本地文件模式 (开发环境)
                if data_path is None:
                    # 尝试多个可能的路径
                    possible_paths = [
                        self.back_dir.parent / 'data' / 'processed' / 'cleaned_energy_data_all.csv',
                        self.back_dir / 'data' / 'processed' / 'cleaned_energy_data_all.csv',
                    ]
                    data_path = None
                    for path in possible_paths:
                        if path.exists():
                            data_path = path
                            break
                    if data_path is None:
                        data_path = possible_paths[0]  # 使用第一个作为默认
                else:
                    data_path = Path(data_path)
            
            # 读取数据
            print(f" 读取数据: {data_path}")
            try:
                df = pd.read_csv(data_path, parse_dates=['Date'])
                print(f"    数据读取成功: {len(df)} 行 × {len(df.columns)} 列")
            except FileNotFoundError:
                raise FileNotFoundError(f"数据文件不存在: {data_path}")
            except Exception as e:
                raise Exception(f"读取数据时出错: {str(e)}")
            
            # ================================================================
            # 【修复】智能特征检查与自动补算
            # 分层检查：硬必需列 vs 可补算的 engineered 特征
            # ================================================================
            
            # 硬必需列：必须存在，无法自动生成
            hard_required = ['Date', 'Site_Load', 'Temperature']
            missing_hard = [col for col in hard_required if col not in df.columns]
            if missing_hard:
                raise ValueError(f"数据缺少核心必需列（无法自动生成）: {missing_hard}")
            
            # 可补算的基础时间特征
            time_features = ['Hour', 'DayOfWeek', 'Price']
            missing_time = [col for col in time_features if col not in df.columns]
            
            # 可补算的 Lag/Rolling 特征
            engineered_features = ['Lag_1h', 'Lag_24h', 'Lag_168h', 
                                   'Rolling_Mean_6h', 'Rolling_Std_6h', 'Rolling_Mean_24h',
                                   'Temp_x_Hour', 'Lag24_x_DayOfWeek']
            missing_engineered = [col for col in engineered_features if col not in df.columns]
            
            # 检查增强特征是否缺失
            missing_enhanced_time = []
            if use_enhanced_features:
                 enhanced_time_cols = ['Month', 'Season', 'IsWeekend', 'IsHoliday']
                 missing_enhanced_time = [c for c in enhanced_time_cols if c not in df.columns]

            # 如果有缺失的可补算特征，自动调用 EnergyDataProcessor 补算
            if missing_time or missing_engineered or missing_enhanced_time:
                print(f"\n 检测到缺失特征，自动补算中...")
                if missing_time:
                    print(f"   - 缺失时间特征: {missing_time}")
                if missing_engineered:
                    print(f"   - 缺失工程特征: {missing_engineered}")
                
                from services.data_processor import EnergyDataProcessor
                processor = EnergyDataProcessor()
                
                # 确保 Date 是 datetime 类型
                if df['Date'].dtype == 'object':
                    df['Date'] = pd.to_datetime(df['Date'])
                
                # 补算时间特征
                if missing_time:
                    if 'Hour' not in df.columns:
                        df['Hour'] = df['Date'].dt.hour
                    if 'DayOfWeek' not in df.columns:
                        df['DayOfWeek'] = df['Date'].dt.dayofweek
                    if 'Price' not in df.columns:
                        df = processor.add_price_feature(df)
                
                # 补算增强时间特征（如果启用）
                if use_enhanced_features:
                    enhanced_time_cols = ['Month', 'Season', 'IsWeekend', 'IsHoliday']
                    missing_enhanced_time = [c for c in enhanced_time_cols if c not in df.columns]
                    if missing_enhanced_time:
                        df = processor.add_enhanced_time_features(df)
                        print(f"    已补算增强时间特征")
                
                # 补算 Lag/Rolling/Interaction 特征
                if missing_engineered:
                    df = processor.add_advanced_features(df, dropna=False, use_enhanced=use_enhanced_features)
                    print(f"    已补算 Lag/Rolling/Interaction 特征")
                
                print(f"    特征补算完成，当前列数: {len(df.columns)}")
            
            # 处理缺失值
            print(f"\n 检查数据质量...")
            # 检查核心列的缺失值
            core_cols = hard_required + time_features
            null_counts = df[core_cols].isnull().sum()
            if null_counts.sum() > 0:
                print(f"     发现缺失值:")
                for col, count in null_counts[null_counts > 0].items():
                    print(f"      - {col}: {count} 个")
                
                # 对于 Temperature，使用均值填充
                if 'Temperature' in null_counts and null_counts['Temperature'] > 0:
                    # 使用线性插值填充温度缺失值（更适合时间序列）
                    df['Temperature'] = df['Temperature'].interpolate(method='linear', limit_direction='both')
                    # 如果仍有缺失（例如开头或结尾），使用均值兜底
                    if df['Temperature'].isnull().sum() > 0:
                         df['Temperature'].fillna(df['Temperature'].mean(), inplace=True)
                    print(f"    Temperature 缺失值已使用线性插值填充")
            else:
                print(f"    无核心列缺失值")
            
            # ================================================================
            # 动态特征选择（检测数据中可用的增强特征）
            # ================================================================
            print(f"\n 检测可用特征...")
            
            # 首先使用基础特征
            available_features = [col for col in self.base_feature_columns if col in df.columns]
            print(f"    基础特征: {len(available_features)}/{len(self.base_feature_columns)}")
            
            # 如果启用增强特征，检测并添加
            enhanced_features_used = []
            if use_enhanced_features:
                for col in self.enhanced_feature_columns:
                    if col in df.columns:
                        available_features.append(col)
                        enhanced_features_used.append(col)
                
                if enhanced_features_used:
                    print(f"    增强特征: {len(enhanced_features_used)} 个")
                    print(f"     {enhanced_features_used}")
                else:
                    print(f"     数据中未找到增强特征，使用基础特征")
            
            # 更新实际使用的特征列表
            self.feature_columns = available_features
            print(f"    总计使用 {len(self.feature_columns)} 个特征")
                
            # 清除因特征工程 (Lag/Rolling) 产生的 NaN 行
            # 这些行通常位于数据集头部
            before_drop = len(df)
            df.dropna(inplace=True)
            after_drop = len(df)
            if before_drop != after_drop:
                print(f"     已删除 {before_drop - after_drop} 行包含 NaN 的样本 (Lag/Rolling start-up)")
            
            # 准备特征和目标变量
            print(f"\n 准备训练数据...")
            X = df[self.feature_columns].copy()
            y = df[self.target_column].copy()
            
            print(f"   - 特征列: {self.feature_columns}")
            print(f"   - 目标变量: {self.target_column}")
            print(f"   - 数据形状: X={X.shape}, y={y.shape}")
            
            # 划分训练集和测试集
            # 【修复】时间序列数据必须按时间顺序拆分，避免数据泄漏
            print(f"\n  划分数据集 (训练集: {int((1-test_size)*100)}%, 测试集: {int(test_size*100)}%)...")
            
            if use_time_series_cv:
                # 时间序列模式：按时间顺序拆分，测试集为严格未来数据
                # 确保数据按 Date 排序（如果 Date 列存在）
                if 'Date' in df.columns:
                    sort_idx = df['Date'].argsort()
                    X = X.iloc[sort_idx].reset_index(drop=True)
                    y = y.iloc[sort_idx].reset_index(drop=True)
                    print(f"    已按时间顺序排序数据")
                
                # 按时间顺序拆分：最后 test_size 比例作为严格未来 holdout
                split_idx = int(len(X) * (1 - test_size))
                X_train = X.iloc[:split_idx].copy()
                X_test = X.iloc[split_idx:].copy()
                y_train = y.iloc[:split_idx].copy()
                y_test = y.iloc[split_idx:].copy()
                print(f"    时间序列拆分：测试集为严格未来数据 (避免数据泄漏)")
            else:
                # 非时间序列模式：保持原有随机拆分行为（向后兼容）
                X_train, X_test, y_train, y_test = train_test_split(
                    X, y, test_size=test_size, random_state=random_state
                )
                print(f"     随机拆分模式 (仅用于非时间序列数据)")
            
            print(f"   - 训练集: {X_train.shape[0]} 样本")
            print(f"   - 测试集: {X_test.shape[0]} 样本")
            
            # 记录配置状态
            self.use_log_transform = use_log_transform
            
            # ================================================================
            # 异常值过滤 (Anomaly Detection) - 仅针对训练集
            # ================================================================
            if remove_outliers:
                print(f"\n 执行异常值过滤 (IQR)...")
                # 计算 IQR
                Q1 = y_train.quantile(0.25)
                Q3 = y_train.quantile(0.75)
                IQR = Q3 - Q1
                lower_bound = Q1 - 1.5 * IQR
                upper_bound = Q3 + 1.5 * IQR
                
                # 过滤索引
                # 注意：这里我们同时需要过滤 X_train 和 y_train
                mask = (y_train >= lower_bound) & (y_train <= upper_bound)
                outliers_count = len(y_train) - mask.sum()
                
                if outliers_count > 0:
                     X_train = X_train[mask]
                     y_train = y_train[mask]
                     print(f"    已移除 {outliers_count} 个异常样本")
                     print(f"     阈值区间: [{lower_bound:.2f}, {upper_bound:.2f}]")
                else:
                     print(f"    未发现显著异常值")
            
            # ================================================================
            # 目标变量对数变换 (Log Transform)
            # ================================================================
            if use_log_transform:
                print(f"\n执行 Log1p 变换...")
                # 检查负值
                if (y_train < 0).any():
                     print(f"     y_train 包含负值，无法进行 Log 变换，自动禁用")
                     self.use_log_transform = False
                else:
                     y_train = np.log1p(y_train)
                     # y_test 不变换，保持原始比例用于评估
                     print(f"    y_train 已进行 np.log1p 变换")

            # ================================================================
            # 自动模型选择与训练
            # ================================================================
            if auto_select_model and model_type == 'auto':
                print(f"\n 自动模型选择模式...")
                best_model, best_params, selection_info = self._auto_select_best_model(
                    X_train, y_train, X_test, y_test, random_state,
                    use_time_series_cv=use_time_series_cv,
                    n_splits=cv_folds,
                    perform_tuning=tune_hyperparameters
                )
                self.model = best_model
                selected_model_type = selection_info['winner']
                hyperparameters = best_params
            else:
                # 手动模式：使用指定的模型类型
                print(f"\n 训练 {model_type} 模型...")
                self.model, hyperparameters = self._create_model(
                    model_type, n_estimators, random_state
                )
                self.model.fit(X_train, y_train)
                selected_model_type = model_type
                selection_info = {'winner': model_type, 'candidates_evaluated': [model_type]}
            
            print(f"    模型训练完成! (类型: {selected_model_type})")
            
            # 评估模型
            print(f"\n 评估模型性能...")
            
            # 训练集预测
            y_train_pred = self.model.predict(X_train)
            
            # 还原
            if self.use_log_transform:
                y_train = np.expm1(y_train)
                y_train_pred = np.expm1(y_train_pred)
                
            train_mae = mean_absolute_error(y_train, y_train_pred)
            train_rmse = np.sqrt(mean_squared_error(y_train, y_train_pred))
            
            # 测试集预测
            y_test_pred = self.model.predict(X_test)
            
            if self.use_log_transform:
                y_test_pred = np.expm1(y_test_pred)
                
            test_mae = mean_absolute_error(y_test, y_test_pred)
            test_rmse = np.sqrt(mean_squared_error(y_test, y_test_pred))
            
            # 计算 R² Score (测试集)
            test_r2 = r2_score(y_test, y_test_pred)
            
            # 计算 MAPE (测试集) - Mean Absolute Percentage Error
            # 避免分母为 0
            mask = y_test.values != 0
            if np.sum(mask) > 0:
                test_mape = np.mean(np.abs((y_test.values[mask] - y_test_pred[mask]) / y_test.values[mask])) * 100
            else:
                test_mape = 0.0
            
            print(f"\n   训练集性能:")
            print(f"      - MAE:  {train_mae:.2f} kW")
            print(f"      - RMSE: {train_rmse:.2f} kW")
            
            print(f"\n   测试集性能:")
            print(f"      - MAE:  {test_mae:.2f} kW")
            print(f"      - RMSE: {test_rmse:.2f} kW")
            print(f"      - R²:   {test_r2:.4f}")
            print(f"      - MAPE: {test_mape:.2f}%")

            # ------------------------------------------------------------
            # 验证信息汇总 (用于元数据和前端展示)
            # ------------------------------------------------------------
            validation_summary = {
                'method': 'TimeSeriesSplit' if use_time_series_cv else 'HoldOut',
                'cv_folds': cv_folds if use_time_series_cv else None,
                'holdout_mae': float(test_mae),
                'holdout_rmse': float(test_rmse),
                'holdout_r2': float(test_r2),
                'holdout_mape': float(test_mape / 100),  # 统一为小数
            }

            # 如果有交叉验证详情，记录当前胜出模型的均值和波动
            if auto_select_model and model_type == 'auto' and selection_info:
                winner_name = selection_info.get('winner')
                cv_detail = (selection_info.get('cv_details') or {}).get(winner_name, {})
                if cv_detail:
                    validation_summary.update({
                        'cv_mae_mean': cv_detail.get('cv_mae_mean'),
                        'cv_mae_std': cv_detail.get('cv_mae_std'),
                        'cv_scores': cv_detail.get('cv_scores')
                    })

            # 数据覆盖范围
            data_coverage = None
            if 'Date' in df.columns:
                date_min = df['Date'].min()
                date_max = df['Date'].max()
                span_days = None
                try:
                    span_days = int((date_max - date_min).days)
                except Exception:
                    span_days = None

                data_coverage = {
                    'start': date_min.isoformat() if hasattr(date_min, 'isoformat') else str(date_min),
                    'end': date_max.isoformat() if hasattr(date_max, 'isoformat') else str(date_max),
                    'span_days': span_days,
                    'rows': int(len(df))
                }
            
            # 特征重要性
            print(f"\n 特征重要性:")
            importances = None
            
            if hasattr(self.model, 'feature_importances_'):
                importances = self.model.feature_importances_
            elif isinstance(self.model, VotingRegressor):
                # 尝试计算平均特征重要性
                try:
                    all_importances = [est.feature_importances_ for est in self.model.estimators_]
                    importances = np.mean(all_importances, axis=0)
                except Exception:
                     pass
            
            if importances is not None:
                feature_importance = pd.DataFrame({
                    'Feature': self.feature_columns,
                    'Importance': importances
                }).sort_values('Importance', ascending=False)
                
                for _, row in feature_importance.iterrows():
                    print(f"      - {row['Feature']}: {row['Importance']:.4f}")
            else:
                print("      (当前模型不支持特征重要性输出)")
                feature_importance = pd.DataFrame(columns=['Feature', 'Importance'])
            
            # 保存模型到 Firebase Storage
            print(f"\n 保存模型到 Firebase Storage: {self.firebase_model_path}")
            temp_model_path = None
            try:
                # Step A: 创建临时文件
                with tempfile.NamedTemporaryFile(mode='wb', suffix='.joblib', delete=False) as tmp_file:
                    temp_model_path = tmp_file.name
                    print(f"   - 临时文件: {temp_model_path}")
                
                # Step B: 保存模型到临时文件
                joblib.dump(self.model, temp_model_path)
                print(f"    模型已保存到临时文件")
                
                # Step B-2: 保存模型到本地持久化路径 (用于开发环境调试)
                # 仅在非 GAE 环境下执行，避免 Read-only file system 错误
                if not self._is_gae_environment():
                    try:
                        # 确保存储目录存在
                        self.local_model_path.parent.mkdir(parents=True, exist_ok=True)
                        joblib.dump(self.model, self.local_model_path)
                        print(f"    模型已备份到本地路径: {self.local_model_path}")
                    except Exception as local_e:
                        print(f"     无法保存本地模型副本: {str(local_e)}")
                else:
                    print(f"   ℹ  GAE 环境：跳过本地模型备份")

                
                # Step C: 上传到 Firebase Storage
                with open(temp_model_path, 'rb') as f:
                    self.storage_service.upload_file(
                        file_data=f,
                        destination_path=self.firebase_model_path,
                        content_type='application/octet-stream'
                    )
                print(f"    模型已上传到 Firebase Storage")
                
            except Exception as e:
                print(f"    模型保存失败: {str(e)}")
                raise
            finally:
                # Step D: 清理临时文件
                if temp_model_path and os.path.exists(temp_model_path):
                    try:
                        os.remove(temp_model_path)
                        print(f"    已清理临时模型文件")
                    except Exception as e:
                        print(f"     清理临时模型文件失败: {str(e)}")
            
            print("\n" + "="*80)
            print(" 模型训练完成!")
            print("="*80 + "\n")
            
            # 保存模型元数据到 Firebase Storage (全局元数据)
            # 获取模型类型名称
            model_type_name = self._get_model_type_name()
            
            try:
                metadata = {
                    'model_type': model_type_name,
                    'model_version': datetime.now().strftime('%Y%m%d_%H%M%S'),
                    'trained_at': datetime.now().isoformat(),
                    'metrics': {
                        'train_mae': float(train_mae),
                        'train_rmse': float(train_rmse),
                        'test_mae': float(test_mae),
                        'test_rmse': float(test_rmse),
                        'r2_score': float(test_r2),
                        'mape': float(test_mape / 100),  # 存储为小数形式, 5% -> 0.05
                    },
                    'training_samples': len(df),
                    'data_source': 'CAISO Real-Time Stream',
                    'feature_importance': feature_importance.to_dict('records'),
                    'model_path': self.firebase_model_path,
                    'status': 'active',
                    # 数据覆盖信息
                    'data_coverage': data_coverage,
                    # 【新增】单位标注信息（向后兼容，不改变数值）
                    'units': {
                        'target_variable': self.target_column,
                        'target_unit': 'kW',  # Site_Load 的单位
                        'load_unit_multiplier': 1.0,  # 1.0 表示原始值，1000 表示 MW->kW
                        'price_unit': 'CNY/kWh',  # 电价单位
                        'temperature_unit': '°C',  # 温度单位
                        'note': 'CAISO 原始数据为 MW，已转换为 kW 存储'
                    },
                    # 特征工程信息
                    'feature_engineering': {
                        'total_features': len(self.feature_columns),
                        'base_features': len(self.base_feature_columns),
                        'enhanced_features': len(enhanced_features_used) if use_enhanced_features else 0,
                        'enhanced_features_list': enhanced_features_used if use_enhanced_features else [],
                        'use_enhanced': use_enhanced_features
                    },
                    # 验证摘要 (交叉验证 / 留出验证)
                    'validation_summary': validation_summary,
                    # 【新增】训练配置记录（便于复现）
                    'training_config': {
                        'test_size': test_size,
                        'random_state': random_state,
                        'use_time_series_cv': use_time_series_cv,
                        'cv_folds': cv_folds if use_time_series_cv else None,
                        'cv_folds': cv_folds if use_time_series_cv else None,
                        'time_series_split': use_time_series_cv,  # 标记是否使用了时间序列拆分
                        'use_log_transform': use_log_transform,   # 记录是否使用了 Log 变换
                        'remove_outliers': remove_outliers,       # 记录是否使用了异常值过滤
                        'hyperparameter_tuning': tune_hyperparameters # 记录是否启用了超参数搜索
                    }
                }
                
                # 添加自动模型选择信息
                if auto_select_model and model_type == 'auto':
                    metadata['auto_selection'] = {
                        'enabled': True,
                        'candidates_evaluated': selection_info.get('candidates_evaluated', []),
                        'winner': selection_info.get('winner', 'unknown'),
                        'improvement_over_baseline': selection_info.get('improvement', 'N/A'),
                        'all_scores': selection_info.get('all_scores', {}),
                        # 新增：交叉验证信息
                        'validation_method': selection_info.get('validation_method', 'HoldOut'),
                        'cv_folds': selection_info.get('cv_folds'),
                        'cv_details': selection_info.get('cv_details')
                    }
                    metadata['hyperparameters'] = hyperparameters
                
                self._save_model_metadata(metadata)
            except Exception as e:
                print(f"     保存模型元数据失败: {str(e)}")
            
            # 返回评估指标
            return {
                'train_mae': train_mae,
                'train_rmse': train_rmse,
                'test_mae': test_mae,
                'test_rmse': test_rmse,
                'r2_score': test_r2,  # 新增
                'mape': test_mape / 100,  # 新增 (小数形式)
                'feature_importance': feature_importance.to_dict('records'),
                'model_type': model_type_name,
                'hyperparameters': hyperparameters if auto_select_model else {'n_estimators': n_estimators},
                'auto_selection': selection_info if auto_select_model and model_type == 'auto' else None,
                # 新增：特征工程和交叉验证信息
                'feature_engineering': {
                    'total_features': len(self.feature_columns),
                    'enhanced_features_used': enhanced_features_used if use_enhanced_features else []
                },
                'validation': {
                    'method': 'TimeSeriesSplit' if use_time_series_cv else 'HoldOut',
                    'cv_folds': cv_folds if use_time_series_cv else None,
                    'summary': validation_summary
                },
                'data_coverage': data_coverage
            }
        
        finally:
            # 清理临时文件
            if temp_data_path and os.path.exists(temp_data_path):
                try:
                    os.remove(temp_data_path)
                    print(f" 清理临时训练数据文件")
                except Exception as e:
                    print(f"  清理临时文件失败: {str(e)}")
    
    def load_model(self) -> bool:
        """
        加载已保存的模型（优先从 Firebase Storage）
        同时加载模型元数据以恢复正确的特征列表
        
        Returns:
            是否加载成功
            
        Raises:
            FileNotFoundError: 模型文件不存在
            Exception: 加载模型时出错
        """
        print(f" 加载模型...")
        temp_model_path = None
        
        try:
            # Step A: 检查 Firebase Storage 中是否存在模型
            print(f"   - 检查 Firebase Storage: {self.firebase_model_path}")
            
            if self.storage_service.file_exists(self.firebase_model_path):
                print(f"    Firebase Storage 中存在模型")
                
                # Step B: 下载到临时目录
                temp_model_path = self.storage_service.download_to_temp(self.firebase_model_path)
                
                if temp_model_path is None:
                    raise Exception("从 Firebase Storage 下载模型失败")
                
                print(f"    模型已下载到: {temp_model_path}")
                
                # Step C: 从临时文件加载模型
                self.model = joblib.load(temp_model_path)
                print(f"    模型加载成功 (来源: Firebase Storage)")
                
                # Step D: 加载模型元数据以恢复特征列表
                self._load_feature_columns_from_metadata()
                
                return True
            
            else:
                # Step E: Firebase 中没有模型，尝试加载本地兜底模型
                print(f"     Firebase Storage 中无模型，尝试加载本地兜底模型")
                
                if not self.local_model_path.exists():
                    raise FileNotFoundError(
                        f"模型文件不存在:\n"
                        f"  - Firebase Storage: {self.firebase_model_path} (不存在)\n"
                        f"  - 本地路径: {self.local_model_path} (不存在)\n"
                        f"请先调用 train_model() 训练模型。"
                    )
                
                self.model = joblib.load(self.local_model_path)
                print(f"    模型加载成功 (来源: 本地兜底文件)")
                
                # 同样尝试加载元数据
                self._load_feature_columns_from_metadata()
                
                return True
        
        except Exception as e:
            raise Exception(f"加载模型时出错: {str(e)}")
        
        finally:
            # 清理临时文件
            if temp_model_path and os.path.exists(temp_model_path):
                try:
                    os.remove(temp_model_path)
                    print(f"    已清理临时模型文件")
                except Exception as e:
                    print(f"     清理临时模型文件失败: {str(e)}")
    
    def _load_history_context(self, end_time: datetime, window_size: int = 200) -> pd.DataFrame:
        """
        加载用于特征构建的历史数据上下文
        
        Args:
            end_time: 截止时间（不包含）
            window_size: 需要加载的历史窗口大小（小时）
            
        Returns:
            包含 Site_Load, Temperature 等列的历史 DataFrame
        """
        print(f" 加载历史上下文 (截止 {end_time})...")
        
        # 尝试加载全量数据文件
        # 在生产环境中，这应该优化为只从数据库/Storage读取部分数据
        # 但为了保持一致性，我们这里复用训练数据文件
        data_path = None
        
        # 尝试从本地或临时目录查找
        possible_paths = [
            self.back_dir.parent / 'data' / 'processed' / 'cleaned_energy_data_all.csv', # 开发环境
            Path('/tmp/cleaned_energy_data_all.csv'), # 临时目录
            self.back_dir / 'data' / 'processed' / 'cleaned_energy_data_all.csv',
        ]
        
        for path in possible_paths:
            if path.exists():
                data_path = path
                break
                
        if data_path is None:
            # 尝试从 Storage 下载
            print("    本地未找到数据，尝试从 Firebase Storage 下载...")
            try:
                data_path = self.storage_service.download_to_temp('data/processed/cleaned_energy_data_all.csv')
            except Exception as e:
                print(f"     无法下载历史数据: {e}")
                
        if data_path:
            df = pd.read_csv(data_path, parse_dates=['Date'])
            # 筛选截止时间之前的数据
            history = df[df['Date'] < end_time].tail(window_size).copy()
            if len(history) < 168:
                print(f"     历史数据不足 168 小时 (实际: {len(history)}), 特征可能不准确")
            return history
        else:
            print("     无法加载历史上下文，将使用全 0 填充 (仅用于测试/冷启动)")
            # 创建虚拟历史数据
            dates = [end_time - timedelta(hours=i) for i in range(window_size, 0, -1)]
            return pd.DataFrame({
                'Date': dates,
                'Site_Load': [0.0] * window_size,
                'Temperature': [25.0] * window_size,
                'Hour': [d.hour for d in dates],
                'DayOfWeek': [d.weekday() for d in dates]
            })

    def predict_next_24h(
        self,
        start_time: Union[str, datetime],
        temp_forecast_list: Optional[List[float]] = None,
        temp_adjust_delta: float = 0.0
    ) -> List[Dict[str, Union[datetime, float]]]:
        """
        预测未来24小时的能源负载 (递归预测模式)
        
        Args:
            start_time: 开始时间
            temp_forecast_list: 温度预测列表
            temp_adjust_delta: 温度调整值 (用于 What-If 分析)
            
        Returns:
            预测结果列表
        """
        if self.model is None:
            raise ValueError("模型未加载，请先调用 load_model() 或 train_model()")
        
        if isinstance(start_time, str):
            start_time = pd.to_datetime(start_time)
        
        # 验证温度预测列表长度
        if temp_forecast_list is not None and len(temp_forecast_list) != 24:
            raise ValueError(f"temp_forecast_list 长度必须为 24，当前为 {len(temp_forecast_list)}")
            
        print(f"\n 递归预测未来24小时负载 (从 {start_time} 开始)...")
        
        # 1. 加载历史上下文 (用于构建 Lag/Rolling 特征)
        # 我们至少需要过去 168 小时的数据
        history_df = self._load_history_context(start_time, window_size=200)
        
        # 转换为 list 以便高效 append
        history_loads = history_df['Site_Load'].tolist()
        history_temps = history_df['Temperature'].tolist() # 历史温度
        
        # 如果未提供温度预测，使用持久性预测 (昨天的温度)
        if temp_forecast_list is None:
            if len(history_temps) >= 24:
                # 使用过去 24 小时的数据作为基准 (Persistence Forecast)
                print("   ℹ  未提供温度预测，使用过去24小时温度作为基准 (Persistence Forecast)")
                temp_forecast_list = history_temps[-24:]
            else:
                # 历史数据不足，回退到默认值
                print("     历史温度不足，使用默认 25.0°C")
                temp_forecast_list = [25.0] * 24
        
        # 应用温度调整 (What-If Analysis)
        if temp_adjust_delta != 0.0:
            print(f"     应用温度调整: {temp_adjust_delta:+.1f}°C")
            temp_forecast_list = [t + temp_adjust_delta for t in temp_forecast_list]
        
        predictions = []
        prediction_results = []
        
        # 2. 递归预测循环
        current_time = start_time
        
        for i in range(24):
            # A. 特征构建
            hour = current_time.hour
            day_of_week = current_time.weekday()
            temperature = temp_forecast_list[i]
            price = self._get_price(hour)
            
            # 构建高级特征
            # 注意：history_loads 的最后一个元素是 t-1 时刻的负载
            
            # Lag Features
            lag_1h = history_loads[-1] if len(history_loads) >= 1 else 0
            lag_24h = history_loads[-24] if len(history_loads) >= 24 else 0
            lag_168h = history_loads[-168] if len(history_loads) >= 168 else 0
            
            # Rolling Features
            # 取最近 N 个点计算均值/标准差
            # 【修复】使用 ddof=1 与训练时 pandas.rolling().std() 保持一致
            roll_6h_mean = np.mean(history_loads[-6:]) if len(history_loads) >= 6 else lag_1h
            roll_6h_std = np.std(history_loads[-6:], ddof=1) if len(history_loads) >= 6 else 0
            roll_24h_mean = np.mean(history_loads[-24:]) if len(history_loads) >= 24 else lag_1h
            
            # Interaction Features (基础)
            temp_x_hour = temperature * hour
            lag24_x_dow = lag_24h * day_of_week
            
            # 组装基础特征向量
            feature_dict = {
                'Hour': hour,
                'DayOfWeek': day_of_week,
                'Temperature': temperature,
                'Price': price,
                'Lag_1h': lag_1h,
                'Lag_24h': lag_24h,
                'Lag_168h': lag_168h,
                'Rolling_Mean_6h': roll_6h_mean,
                'Rolling_Std_6h': roll_6h_std,
                'Rolling_Mean_24h': roll_24h_mean,
                'Temp_x_Hour': temp_x_hour,
                'Lag24_x_DayOfWeek': lag24_x_dow
            }
            
            # 添加增强特征（如果模型需要）
            # 检查模型是否使用增强特征
            if len(self.feature_columns) > 12:
                # 时间特征
                month = current_time.month
                day_of_month = current_time.day
                week_of_year = current_time.isocalendar()[1]
                
                # 季节 (北半球)
                if month in [3, 4, 5]:
                    season = 0  # 春
                elif month in [6, 7, 8]:
                    season = 1  # 夏
                elif month in [9, 10, 11]:
                    season = 2  # 秋
                else:
                    season = 3  # 冬
                
                # 是否周末
                is_weekend = 1 if day_of_week >= 5 else 0
                
                # 是否节假日（美国加州）
                # 优化: 避免在循环中重复初始化
                if not hasattr(self, '_holidays_cache'):
                    try:
                        import holidays
                        # 使用 subdiv 替代 deprecated 的 state 参数
                        self._holidays_cache = holidays.US(subdiv='CA')
                    except ImportError:
                        self._holidays_cache = None
                    except Exception:
                        # 兼容旧版本 holidays (如不支持 subdiv)
                        try:
                            import holidays
                            self._holidays_cache = holidays.US(state='CA')
                        except:
                            self._holidays_cache = None

                if self._holidays_cache:
                    is_holiday = 1 if current_time.date() in self._holidays_cache else 0
                else:
                    is_holiday = is_weekend  # 简化：周末视为假日
                
                # 增强交互特征
                temp_x_season = temperature * season
                lag24_x_is_weekend = lag_24h * is_weekend
                hour_x_is_holiday = hour * is_holiday
                
                # 周期编码
                month_sin = np.sin(2 * np.pi * month / 12)
                month_cos = np.cos(2 * np.pi * month / 12)
                hour_sin = np.sin(2 * np.pi * hour / 24)
                hour_cos = np.cos(2 * np.pi * hour / 24)
                
                # 添加增强特征
                feature_dict.update({
                    'Month': month,
                    'Season': season,
                    'IsWeekend': is_weekend,
                    'IsHoliday': is_holiday,
                    'DayOfMonth': day_of_month,
                    'WeekOfYear': week_of_year,
                    'Temp_x_Season': temp_x_season,
                    'Lag24_x_IsWeekend': lag24_x_is_weekend,
                    'Hour_x_IsHoliday': hour_x_is_holiday,
                    'Month_Sin': month_sin,
                    'Month_Cos': month_cos,
                    'Hour_Sin': hour_sin,
                    'Hour_Cos': hour_cos
                })
            
            # 确保特征顺序与模型一致
            features = pd.DataFrame([{col: feature_dict[col] for col in self.feature_columns}])
            
            # B. 单步推理
            pred_log = float(self.model.predict(features)[0])
            
            # 如果模型使用了 Log 变换，需要还原
            if hasattr(self, 'use_log_transform') and self.use_log_transform:
                pred_load = float(np.expm1(pred_log))
            else:
                pred_load = pred_log
            
            # 【安全检查】强制约束预测值为非负 (物理常识)
            pred_load = max(0.0, pred_load)
            
            # C. 更新历史序列 (递归关键)
            # 将预测值作为"真实值"加入历史，用于下一步预测
            history_loads.append(pred_load)
            predictions.append(pred_load)
            
            # D. 记录结果
            prediction_results.append({
                'datetime': current_time,
                'predicted_load': pred_load,
                'temperature': temperature,
                'price': price,
                'hour': hour,
                'day_of_week': day_of_week
            })
            
            # 时间步进
            current_time += timedelta(hours=1)
            
        print(f"    递归预测完成")
        return prediction_results

    def predict_single(self, *args, **kwargs):
        """
        单点预测已被递归预测取代，且不仅依赖简单输入。
        为保持接口兼容抛出异常或仅做简单处理。
        """
        raise NotImplementedError("单点预测 (predict_single) 已废弃，请使用 predict_next_24h 进行序列预测。")
    
    def get_feature_importance(self) -> Dict[str, float]:
        """
        获取模型的全局特征重要性
        
        使用随机森林内置的 feature_importances_ 属性，
        反映各特征在整体预测中的平均贡献
        
        Returns:
            特征名到重要性分数的映射字典
            
        Raises:
            ValueError: 模型未加载
        """
        if self.model is None:
            raise ValueError("模型未加载，请先调用 load_model() 或 train_model()")
        
        importances = None
        if hasattr(self.model, 'feature_importances_'):
             importances = self.model.feature_importances_
        elif isinstance(self.model, VotingRegressor):
             try:
                 all_importances = [est.feature_importances_ for est in self.model.estimators_]
                 importances = np.mean(all_importances, axis=0)
             except:
                 pass
                 
        if importances is None:
             # 如果无法获取，返回空字典或默认均匀分布
             return {col: 0.0 for col in self.feature_columns}
             
        # 【修复】归一化特征重要性，确保总和为 1.0
        # XGBoost/LightGBM 在 VotingRegressor 中可能返回未归一化的值 (如 gain/cover)
        total_importance = np.sum(importances)
        if total_importance > 0:
            importances = importances / total_importance
             
        importance_dict = dict(zip(
            self.feature_columns,
            [float(v) for v in importances]
        ))
        
        # 按重要性降序排序
        sorted_importance = dict(sorted(
            importance_dict.items(),
            key=lambda x: x[1],
            reverse=True
        ))
        
        return sorted_importance
    
    def explain_prediction(
        self,
        hour: int,
        day_of_week: int,
        temperature: float,
        price: float = None
    ) -> Dict[str, Any]:
        """
        使用 SHAP 解释单次预测
        
        为给定输入提供详细的特征贡献分析，
        展示每个特征如何影响最终预测值
        
        Args:
            hour: 小时 (0-23)
            day_of_week: 星期几 (0-6)
            temperature: 温度 (°C)
            price: 电价，如果为None则自动根据hour计算
            
        Returns:
            包含以下字段的字典:
            - base_value: 模型基准预测值（训练数据的平均值）
            - predicted_value: 实际预测值
            - feature_contributions: 各特征的贡献值字典
            - interpretation: 人类可读的解释文字
        """
        try:
            import shap
            
            # 检查模型是否已加载
            if self.model is None:
                raise ValueError("模型未加载，请先调用 load_model() 或 train_model()")

            # 构建特征 DataFrame
            # 【修复】如果 price 为 None，使用统一的 _get_price 方法计算（与推理一致）
            if price is None:
                price = self._get_price(hour)
            
            # 尝试从历史数据获取滞后特征，如果没有则使用合理的默认值
            # 默认值基于典型的负载模式
            default_load = 150.0  # 典型平均负载 (kW)
            
            # 构建与训练时相同的特征 DataFrame（包含所有12个特征）
            features = pd.DataFrame({
                'Hour': [hour],
                'DayOfWeek': [day_of_week],
                'Temperature': [temperature],
                'Price': [price],
                'Lag_1h': [default_load],  # 1小时前的负载
                'Lag_24h': [default_load],  # 24小时前的负载
                'Lag_168h': [default_load],  # 168小时(一周)前的负载
                'Rolling_Mean_6h': [default_load],  # 6小时滚动平均
                'Rolling_Std_6h': [default_load * 0.1],  # 6小时滚动标准差
                'Rolling_Mean_24h': [default_load],  # 24小时滚动平均
                'Temp_x_Hour': [temperature * hour],  # 温度与小时的交互特征
                'Lag24_x_DayOfWeek': [default_load * day_of_week]  # 24小时滞后与星期的交互
            })
            
            # 确保特征列顺序与训练时一致
            features = features[self.feature_columns]
            
            # 使用 TreeExplainer 解释随机森林模型
            # 只有当 explainer 尚未初始化时才创建，避免重复计算
            if not hasattr(self, '_shap_explainer') or self._shap_explainer is None:
                self._shap_explainer = shap.TreeExplainer(self.model)
                
            # 计算 SHAP 值
            shap_values = self._shap_explainer.shap_values(features)
            
            # 获取期望值 (base value)
            # 对于回归模型，expected_value 应该是一个标量
            base_value = self._shap_explainer.expected_value
            if isinstance(base_value, np.ndarray):
                base_value = base_value[0]
                
            # 获取当前预测的 SHAP 值
            # shap_values 对于回归可能是 (n_samples, n_features)
            # 我们只需要第一个样本
            current_shap_values = shap_values[0]
            
            # 预测值 = base_value + sum(shap_values)
            predicted_value = base_value + np.sum(current_shap_values)
            
            # 构建特征贡献列表
            contributions = []
            for i, col in enumerate(self.feature_columns):
                contributions.append({
                    'feature': col,
                    'value': float(features.iloc[0][col]),
                    'contribution': float(current_shap_values[i])
                })
            
            # 按贡献绝对值排序
            contributions.sort(key=lambda x: abs(x['contribution']), reverse=True)
            
            # 生成人类可读的解释文字
            top_feature = contributions[0]
            direction = "增加" if top_feature['contribution'] > 0 else "减少"
            interpretation = (
                f"{top_feature['feature']} 是影响最大的因素，"
                f"它使得预测负载{direction}了 {abs(top_feature['contribution']):.1f} kW。"
            )

            return {
                'base_value': float(base_value),
                'predicted_value': float(predicted_value),
                'feature_contributions': contributions,
                'interpretation': interpretation
            }
            
        except Exception as e:
            print(f"解释预测失败: {str(e)}")
            return None

    def evaluate_recent_performance(self, hours: int = 24) -> Dict[str, Union[float, str]]:
        """
        评估模型在最近一段时间的表现 (在线监控)
        通过回测最近的真实数据来计算指标
        
        Args:
            hours: 回测的小时数 (默认 24)
            
        Returns:
            包含评估指标的字典 (mape, r2, last_update_time)
        """
        print(f"\n 开始在线模型评估 (最近 {hours} 小时)...")
        
        try:
            # 1. 动态下载最新数据 (从 Firestore/Storage 持久化的 CSV)
            from services.storage_service import StorageService
            storage_service = StorageService()
            
            # 尝试下载最新的 cleaned_energy_data_all.csv
            data_path = storage_service.download_to_temp('data/processed/cleaned_energy_data_all.csv')
            
            if not data_path:
                print("    无法下载数据文件，跳过评估")
                return {'status': 'no_data'}
                
            # 2. 读取数据
            df = pd.read_csv(data_path, parse_dates=['Date'])
            
            # 3. 基于时间截取最近 N 小时的数据 (避免数据中断导致 tail(N) 跨度过大)
            last_time = df['Date'].max()
            start_time = last_time - timedelta(hours=hours)
            
            # 筛选时间窗口内的数据
            recent_df = df[df['Date'] > start_time].copy()
            
            # 检查样本量
            # 理论上应该有 hours 个样本，允许有一点缺失 (e.g. > 50%)
            min_samples = int(hours * 0.5) 
            if len(recent_df) < min_samples:
                print(f"    最近 {hours} 小时内数据样本不足 ({len(recent_df)} < {min_samples})，跳过评估")
                return {
                    'status': 'insufficient_data',
                    'message': f'Insufficient data in last {hours}h: found {len(recent_df)} samples'
                }
                
            # 4. 准备特征和真实值
            # 【修复】检查并自动补算缺失特征
            missing_features = [col for col in self.feature_columns if col not in recent_df.columns]
            if missing_features:
                print(f"    检测到缺失特征: {missing_features}")
                print(f"    尝试自动补算...")
                
                try:
                    from services.data_processor import EnergyDataProcessor
                    processor = EnergyDataProcessor()
                    
                    # 确保 Date 是 datetime 类型
                    if recent_df['Date'].dtype == 'object':
                        recent_df['Date'] = pd.to_datetime(recent_df['Date'])
                    
                    # 补算基础时间特征
                    if 'Hour' not in recent_df.columns:
                        recent_df['Hour'] = recent_df['Date'].dt.hour
                    if 'DayOfWeek' not in recent_df.columns:
                        recent_df['DayOfWeek'] = recent_df['Date'].dt.dayofweek
                    if 'Price' not in recent_df.columns:
                        recent_df = processor.add_price_feature(recent_df)
                    
                    # 补算增强时间特征
                    enhanced_time_cols = ['Month', 'Season', 'IsWeekend', 'IsHoliday']
                    if any(c in missing_features for c in enhanced_time_cols):
                        recent_df = processor.add_enhanced_time_features(recent_df)
                    
                    # 补算 Lag/Rolling 特征
                    lag_cols = ['Lag_1h', 'Lag_24h', 'Lag_168h', 'Rolling_Mean_6h', 'Rolling_Std_6h']
                    if any(c in missing_features for c in lag_cols):
                        recent_df = processor.add_advanced_features(recent_df, dropna=False, use_enhanced=True)
                    
                    # 再次检查是否还有缺失
                    still_missing = [col for col in self.feature_columns if col not in recent_df.columns]
                    if still_missing:
                        print(f"    仍有无法补算的特征: {still_missing}")
                        return {
                            'status': 'feature_mismatch',
                            'message': f'Cannot compute features: {still_missing}'
                        }
                    print(f"    特征补算完成")
                except Exception as fe:
                    print(f"    特征补算失败: {str(fe)}")
                    return {
                        'status': 'feature_error',
                        'message': f'Feature computation failed: {str(fe)}'
                    }
            
            # 删除因 Lag/Rolling 产生的 NaN 行
            recent_df = recent_df.dropna(subset=self.feature_columns)
            if len(recent_df) < min_samples:
                print(f"    删除 NaN 后样本不足 ({len(recent_df)} < {min_samples})")
                return {
                    'status': 'insufficient_data',
                    'message': f'Insufficient valid samples after feature computation'
                }
            
            X_recent = recent_df[self.feature_columns]
            y_true = recent_df[self.target_column].values
            
            # 5. 进行预测
            y_pred = self.model.predict(X_recent)
            
            # 如果模型使用了 Log 变换，需要还原
            if hasattr(self, 'use_log_transform') and self.use_log_transform:
                 y_pred = np.expm1(y_pred)
            
            # 6. 计算指标
            # MAPE: Mean Absolute Percentage Error
            # 避免分母为 0
            mask = y_true != 0
            if np.sum(mask) == 0:
                print("    所有真实负载均为 0，无法计算 MAPE")
                mape = 0.0
            else:
                mape = (np.mean(np.abs((y_true[mask] - y_pred[mask]) / y_true[mask])) * 100)
            
            # R2 Score
            if len(y_true) < 2:
                r2 = 0.0  # 样本太少
            else:
                r2 = r2_score(y_true, y_pred)
            
            # 7. 格式化结果
            # 注意：mape 存储为小数形式 (0.05 = 5%)，以便与前端 percent indicator 直接兼容
            metrics = {
                'status': 'success',
                'mape': round(mape / 100, 4),  # 转换为小数形式 (5.25% -> 0.0525)
                'r2_score': round(r2, 3),  # 使用正确的 key 名称匹配前端
                'sample_count': len(y_true),
                'last_data_point': last_time.strftime('%Y-%m-%d %H:%M:%S')
            }
            
            print(f"    评估完成: MAPE={mape:.2f}%, R2={metrics['r2_score']}")
            return metrics
            
        except Exception as e:
            print(f"    在线评估失败: {str(e)}")
            return {'status': 'error', 'message': str(e)}
\end{lstlisting}

\subsection{优化调度服务 (back/services/optimization\_service.py)}
\begin{lstlisting}[style=codestyle, language=Python, caption={Gurobi 优化调度服务}]
"""
能源优化服务模块 - 电池储能系统优化调度
Energy Optimization Service - Battery Energy Storage System Scheduling

使用 Gurobi 求解器进行混合整数规划 (MIP) 优化
"""

import os
import numpy as np
import pandas as pd
from typing import List, Dict, Optional, Tuple
import warnings

from services.secrets import get_secret

warnings.filterwarnings('ignore')

try:
    import gurobipy as gp
    from gurobipy import GRB
    GUROBI_AVAILABLE = True
except ImportError:
    gp = None  # 设置为 None 以避免 NameError
    GRB = None
    GUROBI_AVAILABLE = False
    print("  警告: gurobipy 未安装，优化功能将不可用")


class EnergyOptimizer:
    """
    能源优化器
    
    使用混合整数规划 (MIP) 优化电池储能系统的充放电调度
    目标: 最小化总购电成本
    """
    
    def __init__(
        self,
        battery_capacity: float = 60.0,
        max_power: float = 20.0,
        efficiency: float = 0.95
    ):
        """
        初始化优化器
        
        Args:
            battery_capacity: 电池容量 (kWh)，默认 60.0 (工业级储能)
            max_power: 最大充放电功率 (kW)，默认 20.0
            efficiency: 充放电效率，默认 0.95 (95%)
        """
        if not GUROBI_AVAILABLE:
            raise ImportError(
                "gurobipy 未安装。请运行: pip install gurobipy"
            )
        
        # 电池参数
        self.battery_capacity = battery_capacity
        self.max_power = max_power
        self.efficiency = efficiency
        
        # Gurobi 环境
        self.env = None
        
        print(f" 电池参数:")
        print(f"   - 容量: {battery_capacity} kWh")
        print(f"   - 最大功率: {max_power} kW")
        print(f"   - 效率: {efficiency * 100}%")
    
    def _create_gurobi_env(self) -> 'gp.Env':
        """
        创建 Gurobi 环境
        
        支持 WLS (Web License Service) 和本地许可证
        
        Returns:
            Gurobi 环境对象
            
        Raises:
            Exception: 许可证错误
        """
        try:
            # 检查 WLS 环境变量（使用 secrets 模块统一管理）
            wls_access_id = get_secret('GRB_WLSACCESSID')
            wls_secret = get_secret('GRB_WLSSECRET')
            wls_license_id = get_secret('GRB_LICENSEID')
            
            if wls_access_id and wls_secret:
                print(" 使用 WLS (Web License Service) 许可证")
                
                # 创建 WLS 环境
                env = gp.Env(empty=True)
                env.setParam('WLSACCESSID', wls_access_id)
                env.setParam('WLSSECRET', wls_secret)
                
                if wls_license_id:
                    env.setParam('LICENSEID', int(wls_license_id))
                
                env.start()
                print("    WLS 许可证验证成功")
                
            else:
                print(" 使用本地许可证 (或 Size-Limited Trial)")
                
                # 使用默认环境
                env = gp.Env()
                print("    许可证验证成功")
            
            return env
            
        except gp.GurobiError as e:
            error_msg = str(e)
            
            if "No Gurobi license found" in error_msg or "license" in error_msg.lower():
                raise Exception(
                    " Gurobi 许可证未找到或无效\n"
                    "   请执行以下步骤之一:\n"
                    "   1. 申请免费学术许可证: https://www.gurobi.com/academia/\n"
                    "   2. 使用 Size-Limited Trial (自动激活，限制 2000 变量)\n"
                    "   3. 配置 WLS 许可证环境变量:\n"
                    "      export GRB_WLSACCESSID=your_access_id\n"
                    "      export GRB_WLSSECRET=your_secret\n"
                    f"   原始错误: {error_msg}"
                )
            else:
                raise Exception(f"Gurobi 环境创建失败: {error_msg}")
    
    def optimize_schedule(
        self,
        load_profile: List[float],
        price_profile: List[float],
        initial_soc: float = 0.5
    ) -> Dict:
        """
        优化电池充放电调度
        
        Args:
            load_profile: 未来24小时的负载预测 (kW)
            price_profile: 未来24小时的电价 (元/kWh)
            initial_soc: 初始电池电量百分比 (0.0-1.0)
            
        Returns:
            优化结果字典，包含:
                - status: 求解状态
                - schedule: 详细调度计划
                - total_cost_without_battery: 无电池时的总成本
                - total_cost_with_battery: 有电池时的总成本
                - savings: 节省金额
                
        Raises:
            ValueError: 输入参数错误
            Exception: 优化失败
        """
        print("\n" + "="*80)
        print(" 开始优化电池调度")
        print("="*80 + "\n")
        
        # 验证输入
        if len(load_profile) != 24:
            raise ValueError(f"负载数据长度必须为24，当前为 {len(load_profile)}")
        
        if len(price_profile) != 24:
            raise ValueError(f"电价数据长度必须为24，当前为 {len(price_profile)}")
        
        if not 0 <= initial_soc <= 1:
            raise ValueError(f"初始SOC必须在 [0, 1] 范围内，当前为 {initial_soc}")
        
        # 转换为 numpy 数组
        load = np.array(load_profile)
        price = np.array(price_profile)
        
        print(f" 输入数据:")
        print(f"   - 负载范围: {load.min():.2f} - {load.max():.2f} kW")
        print(f"   - 电价范围: {price.min():.2f} - {price.max():.2f} 元/kWh")
        print(f"   - 初始 SOC: {initial_soc * 100:.1f}%")
        
        # 计算无电池时的总成本
        cost_without_battery = np.sum(load * price)
        print(f"   - 无电池总成本: {cost_without_battery:.2f} 元")
        
        try:
            # 创建 Gurobi 环境
            if self.env is None:
                self.env = self._create_gurobi_env()
            
            # 创建模型
            print(f"\n  构建优化模型...")
            model = gp.Model("BatteryScheduling", env=self.env)
            model.setParam('OutputFlag', 0)  # 关闭求解器输出
            
            T = 24  # 时间步数
            
            # 决策变量
            print(f"   - 创建决策变量...")
            
            # 充电功率 (kW)
            P_charge = model.addVars(T, lb=0, ub=self.max_power, name="P_charge")
            
            # 放电功率 (kW)
            P_discharge = model.addVars(T, lb=0, ub=self.max_power, name="P_discharge")
            
            # 电池存储电量 (kWh)
            E_stored = model.addVars(T, lb=0, ub=self.battery_capacity, name="E_stored")
            
            # 二进制变量: 是否充电
            Is_charge = model.addVars(T, vtype=GRB.BINARY, name="Is_charge")
            
            # 二进制变量: 是否放电
            Is_discharge = model.addVars(T, vtype=GRB.BINARY, name="Is_discharge")
            
            print(f"    变量数量: {T * 5} 个")
            
            # 约束条件
            print(f"   - 添加约束条件...")
            
            # 1. 状态互斥约束: 不能同时充放电
            for t in range(T):
                model.addConstr(
                    Is_charge[t] + Is_discharge[t] <= 1,
                    name=f"mutex_{t}"
                )
            
            # 2. 功率限制约束
            for t in range(T):
                # 充电功率限制
                model.addConstr(
                    P_charge[t] <= self.max_power * Is_charge[t],
                    name=f"charge_limit_{t}"
                )
                
                # 放电功率限制
                model.addConstr(
                    P_discharge[t] <= self.max_power * Is_discharge[t],
                    name=f"discharge_limit_{t}"
                )
            
            # 3. 能量守恒约束 (电池动态方程)
            # 单位检查 (防止 MW 数据被当作 kW 使用)
            avg_load = float(np.mean(load_profile))
            if avg_load > 10000:
                # logger.error(f"异常巨大的负载数值 ({avg_load:.2f} kW). 疑似单位错误 (MW vs kW).") # Assuming logger is defined
                # 严重警告，但不阻止运行 (可能会得到不合理的优化结果)
                print(f"  警告: 检测到非常大的负载 ({avg_load:.2f})，请确认单位是否为 kW")
                raise ValueError(
                    f"检测到异常巨大的负载数值 ({avg_load:.2f} kW)。"
                    "这可能表示单位错误 (例如，输入的是 MW 而非 kW)。"
                    "请确保负载数据以 kW 为单位。"
                )
            
            # 4. 防逆流约束 (No Grid Export)
            # 确保: grid_power = load + charge - discharge >= 0
            # 即: discharge <= load + charge
            # 用户反馈负值是不合理的，因此我们默认开启"防逆流"模式，禁止向电网卖电
            for t in range(T):
                model.addConstr(
                    load[t] + P_charge[t] - P_discharge[t] >= 0,
                    name=f"no_export_{t}"
                )
            
            initial_energy = initial_soc * self.battery_capacity
            
            for t in range(T):
                if t == 0:
                    # 初始时刻
                    model.addConstr(
                        E_stored[t] == initial_energy + 
                        P_charge[t] * self.efficiency - 
                        P_discharge[t] / self.efficiency,
                        name=f"energy_balance_{t}"
                    )
                else:
                    # 后续时刻
                    model.addConstr(
                        E_stored[t] == E_stored[t-1] + 
                        P_charge[t] * self.efficiency - 
                        P_discharge[t] / self.efficiency,
                        name=f"energy_balance_{t}"
                    )
            
            print(f"    约束数量: {T * 4} 个")
            
            # 目标函数: 最小化总购电成本
            print(f"   - 设置目标函数...")
            
            total_cost = gp.quicksum(
                (load[t] + P_charge[t] - P_discharge[t]) * price[t]
                for t in range(T)
            )
            
            model.setObjective(total_cost, GRB.MINIMIZE)
            print(f"    目标: 最小化总购电成本")
            
            # 求解模型
            print(f"\n 开始求解...")
            model.optimize()
            
            # 检查求解状态
            status = model.status
            
            if status == GRB.OPTIMAL:
                print(f"    求解成功! (状态: OPTIMAL)")
                
                # 提取结果
                schedule = []
                soc_hits_min = 0
                soc_hits_max = 0
                charge_hits = 0
                discharge_hits = 0
                
                for t in range(T):
                    p_charge = P_charge[t].X
                    p_discharge = P_discharge[t].X
                    e_stored = E_stored[t].X
                    soc = e_stored / self.battery_capacity

                    if soc <= 0.10 + 1e-6:
                        soc_hits_min += 1
                    if soc >= 0.90 - 1e-6:
                        soc_hits_max += 1
                    if abs(p_charge - self.max_power) < 1e-5 and p_charge > 0:
                        charge_hits += 1
                    if abs(p_discharge - self.max_power) < 1e-5 and p_discharge > 0:
                        discharge_hits += 1
                    
                    # 记录调度结果 (kW)
                    # grid_power = load + charge - discharge
                    # 正值 = 从电网买电, 负值 = 向电网卖电
                    grid_power = load[t] + p_charge - p_discharge
                    
                    schedule.append({
                        'hour': t,
                        'load': float(load[t]),
                        'price': float(price[t]),
                        'charge_power': float(p_charge),
                        'discharge_power': float(p_discharge),
                        'battery_action': float(p_charge - p_discharge),
                        'soc': float(soc),
                        'stored_energy': float(e_stored),
                        'grid_power': float(grid_power)
                    })
                    
                # 检查是否存在大量反向送电 (Export)
                total_export = sum(max(0, -item['grid_power']) for item in schedule)
                if total_export > 1000: # 如果全天送电超过 1000 kWh (5MW 电池容易跑到这个值)
                    print(f" 注意: 检测到大量反向送电 ({total_export:.2f} kWh)。")
                    print("   如果这不是预期的，请检查电池容量设置是否过大 (MW级别?)")
                
                # 计算总成本
                cost_with_battery = model.objVal
                savings = cost_without_battery - cost_with_battery
                savings_percent = (savings / cost_without_battery) * 100 if cost_without_battery > 0 else 0
                
                print(f"\n 优化结果:")
                print(f"   - 无电池总成本: {cost_without_battery:.2f} 元")
                print(f"   - 有电池总成本: {cost_with_battery:.2f} 元")
                print(f"   - 节省金额: {savings:.2f} 元 ({savings_percent:.1f}%)")

                diagnostics = {
                    'runtime_sec': float(getattr(model, "Runtime", 0.0)),
                    'mip_gap': float(getattr(model, "MIPGap", 0.0)) if hasattr(model, "MIPGap") else None,
                    'node_count': int(getattr(model, "NodeCount", 0)) if hasattr(model, "NodeCount") else None,
                    'iter_count': int(getattr(model, "IterCount", 0)) if hasattr(model, "IterCount") else None,
                }

                constraint_hits = {
                    'soc_min_hits': soc_hits_min,
                    'soc_max_hits': soc_hits_max,
                    'max_charge_hits': charge_hits,
                    'max_discharge_hits': discharge_hits
                }
                
                return {
                    'status': 'Optimal',
                    'schedule': schedule,
                    'total_cost_without_battery': float(cost_without_battery),
                    'total_cost_with_battery': float(cost_with_battery),
                    'savings': float(savings),
                    'savings_percent': float(savings_percent),
                    'diagnostics': diagnostics,
                    'constraint_hits': constraint_hits
                }
                
            elif status == GRB.INFEASIBLE:
                print(f"    模型不可行 (INFEASIBLE)")
                return {
                    'status': 'Infeasible',
                    'error': '模型约束不可行，请检查输入参数'
                }
                
            elif status == GRB.UNBOUNDED:
                print(f"    模型无界 (UNBOUNDED)")
                return {
                    'status': 'Unbounded',
                    'error': '模型目标函数无界'
                }
                
            else:
                print(f"     求解未完成 (状态码: {status})")
                return {
                    'status': 'Unknown',
                    'error': f'求解状态未知 (状态码: {status})'
                }
                
        except gp.GurobiError as e:
            error_msg = str(e)
            
            if "license" in error_msg.lower():
                print(f"\n Gurobi 许可证错误")
                return {
                    'status': 'Error',
                    'error': 'Optimization failed: Gurobi license not found'
                }
            else:
                print(f"\n Gurobi 错误: {error_msg}")
                return {
                    'status': 'Error',
                    'error': f'Gurobi error: {error_msg}'
                }
                
        except Exception as e:
            print(f"\n 优化失败: {str(e)}")
            return {
                'status': 'Error',
                'error': str(e)
            }
    
    def print_schedule(self, result: Dict):
        """
        打印优化调度结果
        
        Args:
            result: optimize_schedule 返回的结果字典
        """
        if result['status'] != 'Optimal':
            print(f"\n 优化失败: {result.get('error', 'Unknown error')}")
            return
        
        print("\n" + "="*80)
        print(" 优化调度计划")
        print("="*80 + "\n")
        
        schedule = result['schedule']
        
        # 打印表头
        print(f"{'时间':<6} {'负载':<10} {'电价':<10} {'电池动作':<12} {'SOC':<10} {'说明':<15}")
        print("-" * 80)
        
        for item in schedule:
            hour = item['hour']
            load = item['load']
            price = item['price']
            action = item['battery_action']
            soc = item['soc']
            
            # 判断电池动作
            if abs(action) < 0.01:
                action_str = "待机"
                action_val = "0.00 kW"
            elif action > 0:
                action_str = "充电"
                action_val = f"+{action:.2f} kW"
            else:
                action_str = "放电"
                action_val = f"{action:.2f} kW"
            
            # 判断电价时段
            if price <= 0.3:
                period = "谷时"
            elif price <= 0.6:
                period = "平时"
            else:
                period = "峰时"
            
            print(f"{hour:02d}:00  {load:>6.2f} kW  {price:>6.2f}元  {action_val:<12} {soc*100:>5.1f}%  {action_str} ({period})")
        
        # 打印总结
        print("\n" + "-" * 80)
        print(f" 成本分析:")
        print(f"   - 无电池总成本: {result['total_cost_without_battery']:.2f} 元")
        print(f"   - 有电池总成本: {result['total_cost_with_battery']:.2f} 元")
        print(f"   - 节省金额: {result['savings']:.2f} 元 ({result['savings_percent']:.1f}%)")
        print("="*80 + "\n")
    
    def close(self):
        """显式关闭 Gurobi 环境（推荐在使用完毕后调用）"""
        if self.env is not None:
            try:
                self.env.dispose()
                self.env = None
                print(" Gurobi 环境已关闭")
            except Exception as e:
                print(f" 关闭 Gurobi 环境时出错: {e}")
    
    def __enter__(self):
        """上下文管理器入口"""
        return self
    
    def __exit__(self, exc_type, exc_val, exc_tb):
        """上下文管理器出口 - 自动关闭环境"""
        self.close()
        return False
    
    def __del__(self):
        """析构函数: 关闭 Gurobi 环境（兜底）"""
        self.close()
\end{lstlisting}

\section{前端源代码}

\subsection{Dart入口 (front/lib/main.dart)}
\begin{lstlisting}[style=codestyle, language=Dart, caption={Flutter 应用入口}]
import 'package:flutter/material.dart';
import 'package:firebase_core/firebase_core.dart';
import 'firebase_options.dart';
import 'config/app_theme.dart';
import 'screens/login_screen.dart';
import 'services/auth_service.dart';

void main() async {
  WidgetsFlutterBinding.ensureInitialized();
  await Firebase.initializeApp(
    options: DefaultFirebaseOptions.currentPlatform,
  );
  runApp(const MyApp());
}

class MyApp extends StatelessWidget {
  const MyApp({super.key});

  @override
  Widget build(BuildContext context) {
    return MaterialApp(
      title: 'Power Data Analysis',
      theme: AppTheme.lightTheme,
      home: StreamBuilder(
        stream: AuthService().authStateChanges,
        builder: (context, snapshot) {
          if (snapshot.connectionState == ConnectionState.active) {
            final user = snapshot.data;
            return user == null ? const LoginScreen() : const MainNavigation();
          }
          return const CircularProgressIndicator();
        },
      ),
    );
  }
}
\end{lstlisting}

\subsection{UI 主题配置 (front/lib/config/app\_theme.dart)}
\begin{lstlisting}[style=codestyle, language=Dart, caption={Flutter Glassmorphism 主题}]
import 'package:flutter/material.dart';

class AppTheme {
  static ThemeData get lightTheme {
    return ThemeData(
      primarySwatch: Colors.blue,
      scaffoldBackgroundColor: const Color(0xFFF5F7FA), // 浅灰背景
      cardTheme: CardTheme(
        elevation: 8,
        shape: RoundedRectangleBorder(borderRadius: BorderRadius.circular(16)),
        color: Colors.white.withOpacity(0.9), // 磨砂感底色
      ),
      useMaterial3: true,
    );
  }
  
  // Glassmorphism 装饰器
  static BoxDecoration get glassDecoration => BoxDecoration(
    color: Colors.white.withOpacity(0.2),
    borderRadius: BorderRadius.circular(16),
    border: Border.all(color: Colors.white.withOpacity(0.3)),
    boxShadow: [
      BoxShadow(
        color: Colors.black.withOpacity(0.1),
        blurRadius: 10,
        spreadRadius: 2,
      )
    ],
  );
}
\end{lstlisting}

\subsection{数据分析屏幕 (front/lib/screens/analysis\_screen.dart)}
\begin{lstlisting}[style=codestyle, language=Dart, caption={数据分析界面}]
class DataAnalysisScreen extends StatefulWidget {
  @override
  _DataAnalysisScreenState createState() => _DataAnalysisScreenState();
}

class _DataAnalysisScreenState extends State<DataAnalysisScreen> {
  // ... 状态变量 ...

  @override
  Widget build(BuildContext context) {
    return Scaffold(
      body: SingleChildScrollView(
        padding: const EdgeInsets.all(16),
        child: Column(
          children: [
            _buildControlPanel(), // 顶部控制栏
            const SizedBox(height: 20),
            if (_isLoading)
               const CircularProgressIndicator()
            else if (_result != null)
               Column(
                 children: [
                   PowerChartWidget(data: _result!.predictions), // 功率图表
                   const SizedBox(height: 16),
                   FeatureImportanceChart(features: _result!.features), // 特征重要性
                   const SizedBox(height: 16),
                   CorrelationMatrixView(matrix: _result!.correlations), // 相关性矩阵
                 ]
               )
          ],
        ),
      ),
    );
  }
}
\end{lstlisting}

\subsection{建模与优化屏幕 (front/lib/screens/modeling\_screen.dart)}
\begin{lstlisting}[style=codestyle, language=Dart, caption={建模与优化界面}]
/// 能源优化仪表盘 - Glassmorphism 设计
/// 交互式能源调度优化界面
library;

import 'dart:ui';
import 'package:flutter/material.dart';
import 'package:fl_chart/fl_chart.dart';
import 'package:percent_indicator/percent_indicator.dart';
import 'package:intl/intl.dart';
import '../config/app_theme.dart';
import '../services/api_service.dart';
import '../models/optimization_result.dart';
import '../widgets/power_chart_widget.dart';
import '../widgets/soc_chart_widget.dart';
import '../widgets/responsive_wrapper.dart';
import '../widgets/analysis/feature_importance_chart.dart';
import '../utils/responsive_helper.dart';
import '../config/constants.dart';

class ModelingScreen extends StatefulWidget {
  const ModelingScreen({super.key});

  @override
  State<ModelingScreen> createState() => _ModelingScreenState();
}

class _ModelingScreenState extends State<ModelingScreen> {
  // 状态变量
  bool _isLoading = false;
  double _initialSoc = AppConstants.defaultInitialSoc; // 默认 50%
  OptimizationResponse? _result;
  OptimizationResponse? _previousResult; // 用于对比
  String? _errorMessage;
  DateTime? _selectedDate;
  
  //  电池参数 (方案一：交互式优化沙盒)
  // 注意：负载规模约 150-300 kW (微网/商业楼宇级)
  double _batteryCapacity = 500; // kWh (商业储能)
  double _maxPower = 200; // kW
  bool _showAdvancedParams = false; // 是否展开高级参数
  
  //  场景模拟 (方案二)
  String? _selectedScenario; // 'summer', 'winter', 'overtime'
  
  //  What-If 预测 (方案三)
  double _temperatureAdjust = 0.0; // -5 ~ +5 度

  @override
  void initState() {
    super.initState();
    // 默认选择明天
    _selectedDate = DateTime.now().add(const Duration(days: 1));
  }

  /// 执行优化
  Future<void> _runOptimization({bool saveForComparison = true}) async {
    setState(() {
      _isLoading = true;
      _errorMessage = null;
      if (saveForComparison && _result != null) {
        _previousResult = _result; // 保存上次结果用于对比
      }
      _result = null;
    });

    try {
      final result = await ApiService.runOptimization(
        initialSoc: _initialSoc,
        targetDate: _selectedDate,
        batteryCapacity: _batteryCapacity,
        batteryPower: _maxPower,
        temperatureAdjust: _temperatureAdjust,
      );

      if (mounted) {
        setState(() {
          _result = result;
          _isLoading = false;
        });

        if (result.isSuccess) {
          _showSuccessSnackBar('优化完成！节省 ${result.optimization?.summary.savingsFormatted ?? "0"}');
        } else {
          setState(() {
            _errorMessage = result.message ?? result.error ?? '优化失败';
          });
        }
      }
    } catch (e) {
      if (mounted) {
        setState(() {
          _errorMessage = e.toString();
          _isLoading = false;
        });
        _showErrorSnackBar(_errorMessage!);
      }
    }
  }

  void _showSuccessSnackBar(String message) {
    ScaffoldMessenger.of(context).showSnackBar(
      SnackBar(
        content: Row(
          children: [
            const Icon(Icons.check_circle, color: Colors.white),
            const SizedBox(width: 8),
            Expanded(child: Text(message)),
          ],
        ),
        backgroundColor: AppColors.success,
        behavior: SnackBarBehavior.floating,
        shape: RoundedRectangleBorder(
          borderRadius: BorderRadius.circular(AppDecorations.radiusMd),
        ),
        margin: const EdgeInsets.all(16),
      ),
    );
  }

  void _showErrorSnackBar(String message) {
    ScaffoldMessenger.of(context).showSnackBar(
      SnackBar(
        content: Row(
          children: [
            const Icon(Icons.error_outline, color: Colors.white),
            const SizedBox(width: 8),
            Expanded(child: Text(message)),
          ],
        ),
        backgroundColor: AppColors.error,
        behavior: SnackBarBehavior.floating,
        shape: RoundedRectangleBorder(
          borderRadius: BorderRadius.circular(AppDecorations.radiusMd),
        ),
        margin: const EdgeInsets.all(16),
        duration: const Duration(seconds: 5),
      ),
    );
  }

  @override
  Widget build(BuildContext context) {
    return Scaffold(
      backgroundColor: AppColors.background,
      body: RefreshIndicator(
        onRefresh: () async {
          if (!_isLoading && _initialSoc > 0) {
            await _runOptimization();
          }
        },
        child: ResponsiveWrapper(
          maxWidth: ResponsiveHelper.getMaxContentWidth(context),
          child: SingleChildScrollView(
            physics: const AlwaysScrollableScrollPhysics(),
            padding: ResponsiveHelper.getPagePadding(context),
            child: Column(
              crossAxisAlignment: CrossAxisAlignment.start,
              children: [
              // 1. 顶部控制卡片
              _buildControlPanel(),
              
              const SizedBox(height: 16),
              
              // 2. 错误提示
              if (_errorMessage != null) _buildErrorCard(),
              
              // 3. 加载指示器
              if (_isLoading) _buildLoadingCard(),
              
              // 4. AI模型健康度卡片 (关键新增！)
              if (_result?.modelInfo != null) ...[
                const SizedBox(height: 16),
                _buildModelHealthCard(_result!.modelInfo!),
              ],
              
              // 5. 关键指标卡片
              if (_result?.optimization != null) ...[
                const SizedBox(height: 16),
                _buildKeyMetrics(_result!.optimization!),

                if (_result!.optimization!.diagnostics != null ||
                    _result!.optimization!.constraintHits != null) ...[
                  const SizedBox(height: 12),
                  _buildSolverDiagnosticsCard(_result!.optimization!),
                ],
                
                const SizedBox(height: 24),
                
                // 5. 电网交互策略图
                PowerChartWidget(
                  chartData: _result!.optimization!.chartData,
                ),
                
                const SizedBox(height: 16),
                
                // 6. 电池电量变化图
                SocChartWidget(
                  chartData: _result!.optimization!.chartData,
                ),
                
                const SizedBox(height: 16),
                
                // 7. 策略详情
                _buildStrategyDetails(_result!.optimization!),
                
                // 8. 模型可解释性 - 特征重要性图表
                if (_result?.modelExplainability != null) ...[
                  const SizedBox(height: 16),
                  ExpansionTile(
                    initiallyExpanded: false,
                    tilePadding: const EdgeInsets.symmetric(horizontal: 16),
                    shape: RoundedRectangleBorder(
                      borderRadius: BorderRadius.circular(12),
                    ),
                    collapsedShape: RoundedRectangleBorder(
                      borderRadius: BorderRadius.circular(12),
                    ),
                    backgroundColor: Colors.white,
                    collapsedBackgroundColor: Colors.white,
                    leading: Icon(Icons.psychology, color: Colors.purple[600]),
                    title: const Text(
                      ' AI 预测解释',
                      style: TextStyle(fontWeight: FontWeight.bold, fontSize: 16),
                    ),
                    subtitle: Text(
                      '了解哪些因素影响了负载预测',
                      style: TextStyle(fontSize: 12, color: Colors.grey[600]),
                    ),
                    children: [
                      Padding(
                        padding: const EdgeInsets.all(16),
                        child: FeatureImportanceChart(
                          featureImportance: _result!.modelExplainability!.featureImportance,
                          featureDescriptions: _result!.modelExplainability!.featureDescriptions,
                          interpretation: _result!.modelExplainability!.interpretation,
                        ),
                      ),
                    ],
                  ),
                ],
              ],
              
              // 空状态提示
              if (_result == null && !_isLoading && _errorMessage == null)
                _buildEmptyState(),
              
              const SizedBox(height: 32),
            ],
          ),
        ),
        ),
      ),
    );
  }

  /// 1. 顶部控制卡片 - 交互式优化沙盒 (Glassmorphism)
  Widget _buildControlPanel() {
    return ClipRRect(
      borderRadius: BorderRadius.circular(AppDecorations.radiusLg),
      child: BackdropFilter(
        filter: ImageFilter.blur(sigmaX: 10, sigmaY: 10),
        child: Container(
          decoration: BoxDecoration(
            color: Colors.white.withOpacity(0.85),
            borderRadius: BorderRadius.circular(AppDecorations.radiusLg),
            border: Border.all(color: Colors.white.withOpacity(0.3)),
            boxShadow: AppDecorations.shadowMd,
          ),
          padding: const EdgeInsets.all(20.0),
          child: Column(
            crossAxisAlignment: CrossAxisAlignment.start,
            children: [
              // 标题
              Row(
                children: [
                  Container(
                    padding: const EdgeInsets.all(8),
                    decoration: BoxDecoration(
                      color: AppColors.primary.withOpacity(0.1),
                      borderRadius: BorderRadius.circular(AppDecorations.radiusMd),
                    ),
                    child: const Icon(Icons.tune_rounded, color: AppColors.primary, size: 20),
                  ),
                  const SizedBox(width: 12),
                  Expanded(
                    child: Text(
                      '优化沙盒',
                      style: AppTextStyles.h3,
                    ),
                  ),
                  // 高级参数展开按钮
                  TextButton.icon(
                    onPressed: () => setState(() => _showAdvancedParams = !_showAdvancedParams),
                    icon: Icon(
                      _showAdvancedParams ? Icons.expand_less : Icons.expand_more,
                      size: 20,
                      color: AppColors.textMuted,
                    ),
                    label: Text(
                      _showAdvancedParams ? '收起' : '高级',
                      style: AppTextStyles.labelMedium,
                    ),
                  ),
                ],
              ),
              const SizedBox(height: 20),
            
              // ========== 场景选择 ==========
              Text('快速场景', style: AppTextStyles.labelMedium.copyWith(color: AppColors.textMuted)),
              const SizedBox(height: 12),
              Wrap(
                spacing: 8,
                runSpacing: 8,
                children: [
                  _buildScenarioChip('summer', '夏季高温', AppColors.warning, Icons.wb_sunny_rounded),
                  _buildScenarioChip('winter', '冬季寒潮', AppColors.info, Icons.ac_unit_rounded),
                  _buildScenarioChip('overtime', '夜间加班', const Color(0xFF8B5CF6), Icons.nightlight_round),
                ],
              ),
              
              Divider(height: 32, color: AppColors.border),
            
            // ========== 初始电量滑块 ==========
            _buildSliderRow(
              icon: Icons.battery_charging_full,
              iconColor: Colors.green,
              label: '初始电量',
              value: _initialSoc,
              min: 0.0,
              max: 1.0,
              divisions: 20,
              displayValue: '${(_initialSoc * 100).toInt()}%',
              onChanged: (v) => setState(() => _initialSoc = v),
            ),
            
            // ========== 高级参数 (可折叠) ==========
            if (_showAdvancedParams) ...[
              const SizedBox(height: 16),
              
              // 电池容量滑块 (商业级储能)
              _buildSliderRow(
                icon: Icons.battery_full,
                iconColor: Colors.blue,
                label: '电池容量 (商业微网)',
                value: _batteryCapacity,
                min: 100,
                max: 2000,
                divisions: 19,
                displayValue: '${_batteryCapacity.toInt()} kWh',
                onChanged: (v) => setState(() => _batteryCapacity = v),
              ),
              
              const SizedBox(height: 16),
              
              // 最大功率滑块 (商业级)
              _buildSliderRow(
                icon: Icons.flash_on,
                iconColor: Colors.amber,
                label: '最大功率 (微网级)',
                value: _maxPower,
                min: 50,
                max: 1000,
                divisions: 19,
                displayValue: '${_maxPower.toInt()} kW',
                onChanged: (v) => setState(() => _maxPower = v),
              ),
              
              const SizedBox(height: 16),
              
              // What-If 温度调整 (方案三)
              _buildSliderRow(
                icon: Icons.thermostat,
                iconColor: Colors.red,
                label: '温度调整 (What-If)',
                value: _temperatureAdjust,
                min: -5.0,
                max: 5.0,
                divisions: 10,
                displayValue: '${_temperatureAdjust >= 0 ? "+" : ""}${_temperatureAdjust.toInt()}°C',
                onChanged: (v) => setState(() => _temperatureAdjust = v),
              ),
            ],
            
            const SizedBox(height: 16),
            
            // 日期选择
            Row(
              children: [
                const Icon(Icons.calendar_today, color: Colors.orange, size: 20),
                const SizedBox(width: 8),
                const Text('目标日期', style: TextStyle(fontSize: 14, fontWeight: FontWeight.w500)),
                const Spacer(),
                InkWell(
                  onTap: _isLoading ? null : () => _selectDate(context),
                  child: Container(
                    padding: const EdgeInsets.symmetric(horizontal: 12, vertical: 8),
                    decoration: BoxDecoration(
                      border: Border.all(color: Colors.grey[300]!),
                      borderRadius: BorderRadius.circular(8),
                    ),
                    child: Row(
                      mainAxisSize: MainAxisSize.min,
                      children: [
                        Text(
                          _selectedDate != null
                              ? DateFormat('MM-dd').format(_selectedDate!)
                              : '选择',
                          style: const TextStyle(fontSize: 14),
                        ),
                        const SizedBox(width: 4),
                        Icon(Icons.arrow_drop_down, color: Colors.grey[600], size: 20),
                      ],
                    ),
                  ),
                ),
              ],
            ),
            
              const SizedBox(height: 24),
            
              // 开始优化按钮 (CTA 橙色)
              SizedBox(
                width: double.infinity,
                height: 52,
                child: ElevatedButton.icon(
                  onPressed: _isLoading ? null : _runOptimization,
                  icon: _isLoading
                      ? const SizedBox(
                          width: 20,
                          height: 20,
                          child: CircularProgressIndicator(strokeWidth: 2.5, color: Colors.white),
                        )
                      : const Icon(Icons.bolt_rounded, size: 22),
                  label: Text(
                    _isLoading ? '优化中...' : '开始智能调度',
                    style: AppTextStyles.button,
                  ),
                  style: ElevatedButton.styleFrom(
                    backgroundColor: AppColors.cta,
                    foregroundColor: Colors.white,
                    shape: RoundedRectangleBorder(
                      borderRadius: BorderRadius.circular(AppDecorations.radiusMd),
                    ),
                    elevation: 0,
                  ),
                ),
              ),
            
              // 参数摘要
              if (_showAdvancedParams || _selectedScenario != null) ...[
                const SizedBox(height: 12),
                Container(
                  padding: const EdgeInsets.all(10),
                  decoration: BoxDecoration(
                    color: AppColors.surfaceVariant,
                    borderRadius: BorderRadius.circular(AppDecorations.radiusMd),
                  ),
                  child: Row(
                    children: [
                      Icon(Icons.info_outline, color: AppColors.textMuted, size: 16),
                      const SizedBox(width: 8),
                      Expanded(
                        child: Text(
                          '${_batteryCapacity.toInt()}kWh | ${_maxPower.toInt()}kW | ${(_initialSoc * 100).toInt()}%'
                          '${_temperatureAdjust != 0 ? " | ${_temperatureAdjust >= 0 ? "+" : ""}${_temperatureAdjust.toInt()}°C" : ""}',
                          style: AppTextStyles.bodySmall,
                        ),
                      ),
                    ],
                  ),
                ),
              ],
            ],
          ),
        ),
      ),
    );
  }
  
  /// 场景选择芯片 (无 emoji)
  Widget _buildScenarioChip(String id, String label, Color color, IconData icon) {
    final isSelected = _selectedScenario == id;
    return FilterChip(
      avatar: Icon(
        icon,
        size: 16,
        color: isSelected ? Colors.white : color,
      ),
      label: Text(label, style: TextStyle(
        fontSize: 12,
        fontWeight: isSelected ? FontWeight.w600 : FontWeight.normal,
        color: isSelected ? Colors.white : AppColors.textPrimary,
      )),
      selected: isSelected,
      onSelected: _isLoading ? null : (selected) {
        setState(() {
          _selectedScenario = selected ? id : null;
          // 根据场景预设参数
          if (selected) {
            switch (id) {
              case 'summer':
                _temperatureAdjust = 5.0; // 高温
                break;
              case 'winter':
                _temperatureAdjust = -5.0; // 低温
                break;
              case 'overtime':
                _temperatureAdjust = 0.0;
                break;
            }
            _showAdvancedParams = true; // 展开显示参数变化
          }
        });
      },
      backgroundColor: color.withOpacity(0.1),
      selectedColor: color,
      checkmarkColor: Colors.white,
      shape: RoundedRectangleBorder(
        borderRadius: BorderRadius.circular(AppDecorations.radiusFull),
        side: BorderSide(color: isSelected ? color : color.withOpacity(0.3)),
      ),
      padding: const EdgeInsets.symmetric(horizontal: 8, vertical: 4),
    );
  }
  
  /// 通用滑块行
  Widget _buildSliderRow({
    required IconData icon,
    required Color iconColor,
    required String label,
    required double value,
    required double min,
    required double max,
    required int divisions,
    required String displayValue,
    required ValueChanged<double> onChanged,
  }) {
    return Row(
      children: [
        Icon(icon, color: iconColor, size: 20),
        const SizedBox(width: 8),
        SizedBox(
          width: 100,
          child: Text(label, style: const TextStyle(fontSize: 13, fontWeight: FontWeight.w500)),
        ),
        Expanded(
          child: SliderTheme(
            data: SliderTheme.of(context).copyWith(
              trackHeight: 4,
              thumbShape: const RoundSliderThumbShape(enabledThumbRadius: 8),
            ),
            child: Slider(
              value: value,
              min: min,
              max: max,
              divisions: divisions,
              activeColor: iconColor,
              inactiveColor: iconColor.withOpacity(0.2),
              onChanged: _isLoading ? null : onChanged,
            ),
          ),
        ),
        Container(
          width: 70,
          padding: const EdgeInsets.symmetric(horizontal: 8, vertical: 4),
          decoration: BoxDecoration(
            color: iconColor.withOpacity(0.1),
            borderRadius: BorderRadius.circular(6),
          ),
          child: Text(
            displayValue,
            textAlign: TextAlign.center,
            style: TextStyle(fontSize: 13, fontWeight: FontWeight.bold, color: iconColor),
          ),
        ),
      ],
    );
  }

  /// 选择日期
  Future<void> _selectDate(BuildContext context) async {
    final DateTime? picked = await showDatePicker(
      context: context,
      initialDate: _selectedDate ?? DateTime.now().add(const Duration(days: 1)),
      firstDate: DateTime.now(),
      lastDate: DateTime.now().add(const Duration(days: 7)),
      builder: (context, child) {
        return Theme(
          data: Theme.of(context).copyWith(
            colorScheme: ColorScheme.light(
              primary: Colors.blue[700]!,
              onPrimary: Colors.white,
            ),
          ),
          child: child!,
        );
      },
    );
    
    if (picked != null && picked != _selectedDate) {
      setState(() {
        _selectedDate = picked;
      });
    }
  }

  /// 2. 错误卡片
  Widget _buildErrorCard() {
    return Card(
      color: Colors.red[50],
      elevation: 2,
      shape: RoundedRectangleBorder(
        borderRadius: BorderRadius.circular(12),
        side: BorderSide(color: Colors.red[300]!, width: 1),
      ),
      child: Padding(
        padding: const EdgeInsets.all(16.0),
        child: Row(
          children: [
            Icon(Icons.error_outline, color: Colors.red[700], size: 32),
            const SizedBox(width: 12),
            Expanded(
              child: Column(
                crossAxisAlignment: CrossAxisAlignment.start,
                children: [
                  Text(
                    '优化失败',
                    style: TextStyle(
                      fontSize: 16,
                      fontWeight: FontWeight.bold,
                      color: Colors.red[900],
                    ),
                  ),
                  const SizedBox(height: 4),
                  Text(
                    _errorMessage ?? '未知错误',
                    style: TextStyle(
                      fontSize: 14,
                      color: Colors.red[800],
                    ),
                  ),
                ],
              ),
            ),
            IconButton(
              icon: const Icon(Icons.close),
              color: Colors.red[700],
              onPressed: () {
                setState(() {
                  _errorMessage = null;
                });
              },
            ),
          ],
        ),
      ),
    );
  }

  /// 3. 加载卡片
  Widget _buildLoadingCard() {
    return Card(
      elevation: 2,
      shape: RoundedRectangleBorder(borderRadius: BorderRadius.circular(12)),
      child: Padding(
        padding: const EdgeInsets.all(24.0),
        child: Column(
          children: [
            const CircularProgressIndicator(),
            const SizedBox(height: 16),
            Text(
              '正在执行优化...',
              style: TextStyle(
                fontSize: 16,
                fontWeight: FontWeight.w500,
                color: Colors.grey[700],
              ),
            ),
            const SizedBox(height: 8),
            Text(
              '预计需要 30-60 秒',
              style: TextStyle(
                fontSize: 14,
                color: Colors.grey[600],
              ),
            ),
          ],
        ),
      ),
    );
  }

  /// 4. AI模型健康度卡片 - 核心展示"模型生命力"
  Widget _buildModelHealthCard(ModelInfo modelInfo) {
    return Card(
      elevation: 6,
      shape: RoundedRectangleBorder(
        borderRadius: BorderRadius.circular(16),
        side: BorderSide(color: Colors.purple[200]!, width: 2),
      ),
      child: Container(
        decoration: BoxDecoration(
          borderRadius: BorderRadius.circular(16),
          gradient: LinearGradient(
            begin: Alignment.topLeft,
            end: Alignment.bottomRight,
            colors: [Colors.purple[50]!, Colors.blue[50]!],
          ),
        ),
        child: Padding(
          padding: const EdgeInsets.all(20.0),
          child: Column(
            crossAxisAlignment: CrossAxisAlignment.start,
            children: [
              // 标题行
              Row(
                children: [
                  Container(
                    padding: const EdgeInsets.all(8),
                    decoration: BoxDecoration(
                      color: Colors.purple[600],
                      borderRadius: BorderRadius.circular(8),
                    ),
                    child: const Icon(Icons.psychology, color: Colors.white, size: 24),
                  ),
                  const SizedBox(width: 12),
                  const Expanded(
                    child: Column(
                      crossAxisAlignment: CrossAxisAlignment.start,
                      children: [
                        Text(
                          ' AI 模型状态',
                          style: TextStyle(
                            fontSize: 20,
                            fontWeight: FontWeight.bold,
                          ),
                        ),
                        Text(
                          '机器学习预测引擎（眼睛）',
                          style: TextStyle(
                            fontSize: 12,
                            color: Colors.grey,
                          ),
                        ),
                      ],
                    ),
                  ),
                  // 状态指示器
                  Container(
                    padding: const EdgeInsets.symmetric(horizontal: 12, vertical: 6),
                    decoration: BoxDecoration(
                      color: modelInfo.isValid ? Colors.green : Colors.orange,
                      borderRadius: BorderRadius.circular(12),
                    ),
                    child: Text(
                      modelInfo.isValid ? '运行中' : '待训练',
                      style: const TextStyle(
                        color: Colors.white,
                        fontSize: 12,
                        fontWeight: FontWeight.bold,
                      ),
                    ),
                  ),
                ],
              ),
              
              const SizedBox(height: 20),
              
              // 核心指标展示
              if (modelInfo.metrics != null)
                Container(
                  padding: const EdgeInsets.all(16),
                  decoration: BoxDecoration(
                    color: Colors.white.withOpacity(0.6),
                    borderRadius: BorderRadius.circular(12),
                  ),
                  child: Row(
                    children: [
                      // R2 Score (圆形进度条)
                      Expanded(
                        child: Column(
                          children: [
                            CircularPercentIndicator(
                              radius: 40.0,
                              lineWidth: 8.0,
                              percent: (modelInfo.metrics!.r2Score ?? 0).clamp(0.0, 1.0),
                              center: Text(
                                "${((modelInfo.metrics!.r2Score ?? 0) * 100).toStringAsFixed(0)}%",
                                style: TextStyle(
                                  fontWeight: FontWeight.bold,
                                  color: Colors.purple[700],
                                ),
                              ),
                              progressColor: (modelInfo.metrics!.r2Score ?? 0) > 0.8 
                                  ? Colors.green 
                                  : Colors.orange,
                              backgroundColor: Colors.purple[50]!,
                              circularStrokeCap: CircularStrokeCap.round,
                              animation: true,
                            ),
                            const SizedBox(height: 8),
                            Text(
                              "R² Score",
                              style: TextStyle(
                                color: Colors.grey[700],
                                fontWeight: FontWeight.bold,
                                fontSize: 12,
                              ),
                            ),
                          ],
                        ),
                      ),
                      
                      const SizedBox(width: 16),
                      
                      // MAPE (线性进度条)
                      Expanded(
                        child: Column(
                          crossAxisAlignment: CrossAxisAlignment.start,
                          children: [
                            Row(
                              mainAxisAlignment: MainAxisAlignment.spaceBetween,
                              children: [
                                Text(
                                  "MAPE (误差)",
                                  style: TextStyle(
                                    color: Colors.grey[700],
                                    fontWeight: FontWeight.bold,
                                    fontSize: 12,
                                  ),
                                ),
                                Text(
                                  "${((modelInfo.metrics!.mape ?? 0) * 100).toStringAsFixed(1)}%",
                                  style: TextStyle(
                                    color: (modelInfo.metrics!.mape ?? 0) < 0.1 
                                        ? Colors.green 
                                        : Colors.orange,
                                    fontWeight: FontWeight.bold,
                                  ),
                                ),
                              ],
                            ),
                            const SizedBox(height: 8),
                            LinearPercentIndicator(
                              lineHeight: 8.0,
                              percent: (1.0 - (modelInfo.metrics!.mape ?? 0)).clamp(0.0, 1.0),
                              progressColor: (modelInfo.metrics!.mape ?? 0) < 0.1 
                                  ? Colors.green 
                                  : Colors.orange,
                              backgroundColor: Colors.purple[50]!,
                              barRadius: const Radius.circular(4),
                              animation: true,
                            ),
                            const SizedBox(height: 4),
                            Text(
                              (modelInfo.metrics!.mape ?? 0) < 0.1 ? "精度优良" : "精度一般",
                              style: TextStyle(
                                fontSize: 10,
                                color: Colors.grey[600],
                              ),
                            ),
                          ],
                        ),
                      ),
                    ],
                  ),
                )
              else
                const Center(
                  child: Padding(
                    padding: EdgeInsets.all(16.0),
                    child: Text("暂无详细性能指标"),
                  ),
                ),
              
              const SizedBox(height: 16),
              
              // 模型详情网格
              Row(
                children: [
                  Expanded(
                    child: _buildModelStatItem(
                      Icons.model_training,
                      '模型类型',
                      modelInfo.usedAutoSelection 
                          ? modelInfo.winnerModel 
                          : 'Random Forest',
                      Colors.blue[700]!,
                    ),
                  ),
                  const SizedBox(width: 12),
                  Expanded(
                    child: _buildModelStatItem(
                      Icons.storage,
                      '训练数据',
                      modelInfo.trainingSamples != null
                          ? '${modelInfo.trainingSamples} 样本'
                          : 'N/A',
                      Colors.green[700]!,
                    ),
                  ),
                ],
              ),
              
              const SizedBox(height: 12),
              
              Row(
                children: [
                  Expanded(
                    child: _buildModelStatItem(
                      Icons.schedule,
                      '最近更新',
                      modelInfo.trainedAtFormatted,
                      Colors.orange[700]!,
                    ),
                  ),
                  const SizedBox(width: 12),
                  Expanded(
                    child: _buildModelStatItem(
                      Icons.precision_manufacturing,
                      '预测精度 (MAE)',
                      modelInfo.maeFormatted,
                      Colors.purple[700]!,
                    ),
                  ),
                ],
              ),
              
              // 自动模型选择信息卡片 (新增)
              if (modelInfo.usedAutoSelection) ...[
                const SizedBox(height: 16),
                _buildAutoSelectionCard(modelInfo),
              ],
              
              // 训练配置信息 (新增)
              if (modelInfo.trainingConfig != null) ...[
                const SizedBox(height: 16),
                _buildOptimizationConfigCard(modelInfo),
              ],
              
              // 验证与数据覆盖
              if (modelInfo.validationSummary != null || modelInfo.dataCoverage != null) ...[
                const SizedBox(height: 16),
                _buildValidationSummaryCard(modelInfo),
              ],
              
              const SizedBox(height: 16),
              
              // 数据源说明
              Container(
                padding: const EdgeInsets.all(12),
                decoration: BoxDecoration(
                  color: Colors.blue[50],
                  borderRadius: BorderRadius.circular(8),
                ),
                child: Row(
                  children: [
                    Icon(Icons.cloud_download, color: Colors.blue[700], size: 20),
                    const SizedBox(width: 8),
                    Expanded(
                      child: Column(
                        crossAxisAlignment: CrossAxisAlignment.start,
                        children: [
                          Text(
                            '数据来源: ${modelInfo.dataSource ?? "CAISO 实时流"}',
                            style: TextStyle(
                              fontSize: 13,
                              fontWeight: FontWeight.bold,
                              color: Colors.blue[900],
                            ),
                          ),
                          const SizedBox(height: 4),
                          Text(
                            '模型每日凌晨自动重训，持续学习最新用电模式',
                            style: TextStyle(
                              fontSize: 11,
                              color: Colors.blue[700],
                            ),
                          ),
                        ],
                      ),
                    ),
                  ],
                ),
              ),
            ],
          ),
        ),
      ),
    );
  }

  /// 构建自动模型选择信息卡片
  Widget _buildAutoSelectionCard(ModelInfo modelInfo) {
    final autoSelection = modelInfo.autoSelection!;
    return Container(
      padding: const EdgeInsets.all(12),
      decoration: BoxDecoration(
        color: Colors.green[50],
        borderRadius: BorderRadius.circular(8),
        border: Border.all(color: Colors.green[200]!),
      ),
      child: Column(
        crossAxisAlignment: CrossAxisAlignment.start,
        children: [
          // 标题行
          Row(
            children: [
              Icon(Icons.auto_awesome, color: Colors.green[700], size: 18),
              const SizedBox(width: 8),
              Text(
                ' 自动模型选择',
                style: TextStyle(
                  fontSize: 14,
                  fontWeight: FontWeight.bold,
                  color: Colors.green[900],
                ),
              ),
              const Spacer(),
              Container(
                padding: const EdgeInsets.symmetric(horizontal: 8, vertical: 2),
                decoration: BoxDecoration(
                  color: Colors.green[600],
                  borderRadius: BorderRadius.circular(10),
                ),
                child: Text(
                  '已启用',
                  style: const TextStyle(
                    color: Colors.white,
                    fontSize: 10,
                    fontWeight: FontWeight.bold,
                  ),
                ),
              ),
            ],
          ),
          const SizedBox(height: 12),
          
          // 详情网格
          Row(
            children: [
              Expanded(
                child: _buildAutoSelectionItem(
                  ' 胜出模型',
                  autoSelection.winner,
                  Colors.amber[700]!,
                ),
              ),
              const SizedBox(width: 8),
              Expanded(
                child: _buildAutoSelectionItem(
                  ' 性能提升',
                  autoSelection.improvementOverBaseline,
                  Colors.green[700]!,
                ),
              ),
            ],
          ),
          const SizedBox(height: 8),
          Row(
            children: [
              Expanded(
                child: _buildAutoSelectionItem(
                  ' 验证方法',
                  autoSelection.validationMethodFormatted,
                  Colors.blue[700]!,
                ),
              ),
              const SizedBox(width: 8),
              Expanded(
                child: _buildAutoSelectionItem(
                  ' 候选模型',
                  '${autoSelection.candidatesEvaluated.length} 个',
                  Colors.purple[700]!,
                ),
              ),
            ],
          ),
          
          // 展开查看所有候选模型得分
          if (autoSelection.allScores != null && autoSelection.allScores!.isNotEmpty) ...[
            const SizedBox(height: 12),
            ExpansionTile(
              tilePadding: EdgeInsets.zero,
              childrenPadding: EdgeInsets.zero,
              title: Text(
                '查看所有候选模型得分',
                style: TextStyle(
                  fontSize: 12,
                  color: Colors.green[700],
                  fontWeight: FontWeight.w500,
                ),
              ),
              children: [
                Container(
                  padding: const EdgeInsets.all(8),
                  decoration: BoxDecoration(
                    color: Colors.white,
                    borderRadius: BorderRadius.circular(6),
                  ),
                  child: Column(
                    children: autoSelection.allScores!.entries.map((entry) {
                      final scores = entry.value as Map<String, dynamic>;
                      final mae = scores['mae'] ?? 0.0;
                      final isWinner = entry.key == autoSelection.winner;
                      return Padding(
                        padding: const EdgeInsets.symmetric(vertical: 4),
                        child: Row(
                          children: [
                            if (isWinner)
                              const Text(' ', style: TextStyle(fontSize: 12))
                            else
                              const SizedBox(width: 18),
                            Expanded(
                              child: Text(
                                entry.key,
                                style: TextStyle(
                                  fontSize: 12,
                                  fontWeight: isWinner ? FontWeight.bold : FontWeight.normal,
                                  color: isWinner ? Colors.amber[800] : Colors.grey[700],
                                ),
                              ),
                            ),
                            Text(
                              'MAE: ${mae.toStringAsFixed(2)} kW',
                              style: TextStyle(
                                fontSize: 11,
                                color: isWinner ? Colors.green[700] : Colors.grey[600],
                                fontWeight: isWinner ? FontWeight.bold : FontWeight.normal,
                              ),
                            ),
                          ],
                        ),
                      );
                    }).toList(),
                  ),
                ),
              ],
            ),
          ],
        ],
      ),
    );
  }

  /// 构建优化配置信息卡片
  Widget _buildOptimizationConfigCard(ModelInfo modelInfo) {
    if (modelInfo.trainingConfig == null) return const SizedBox();
    
    final config = modelInfo.trainingConfig!;
    
    return Container(
      padding: const EdgeInsets.all(12),
      decoration: BoxDecoration(
        color: Colors.indigo[50],
        borderRadius: BorderRadius.circular(8),
        border: Border.all(color: Colors.indigo[100]!),
      ),
      child: Column(
        crossAxisAlignment: CrossAxisAlignment.start,
        children: [
          Row(
            children: [
              Icon(Icons.tune, color: Colors.indigo[700], size: 18),
              const SizedBox(width: 8),
              Text(
                ' 训练配置',
                style: TextStyle(
                  fontSize: 14,
                  fontWeight: FontWeight.bold,
                  color: Colors.indigo[900],
                ),
              ),
            ],
          ),
          const SizedBox(height: 12),
          
          Wrap(
            spacing: 8,
            runSpacing: 8,
            children: [
              _buildConfigChip(
                'Log1p 变换', 
                config.useLogTransform ?? false,
                Icons.functions,
              ),
              _buildConfigChip(
                '异常值剔除', 
                config.removeOutliers ?? false,
                Icons.filter_alt,
              ),
              _buildConfigChip(
                '超参数调优', 
                config.tuneHyperparameters ?? false,
                Icons.explore,
              ),
               _buildConfigChip(
                '时序交叉验证', 
                config.useTimeSeriesCV ?? false,
                Icons.timeline,
              ),
            ],
          ),
        ],
      ),
    );
  }

  Widget _buildConfigChip(String label, bool enabled, IconData icon) {
    return Container(
      padding: const EdgeInsets.symmetric(horizontal: 10, vertical: 6),
      decoration: BoxDecoration(
        color: enabled ? Colors.white : Colors.grey[200],
        borderRadius: BorderRadius.circular(16),
        border: Border.all(
          color: enabled ? Colors.indigo[200]! : Colors.grey[300]!,
        ),
        boxShadow: enabled ? [
          BoxShadow(
            color: Colors.indigo.withOpacity(0.1),
            blurRadius: 4,
            offset: const Offset(0, 2),
          )
        ] : null,
      ),
      child: Row(
        mainAxisSize: MainAxisSize.min,
        children: [
          Icon(
            icon, 
            size: 14, 
            color: enabled ? Colors.indigo[600] : Colors.grey[500]
          ),
          const SizedBox(width: 6),
          Text(
            label,
            style: TextStyle(
              fontSize: 12,
              fontWeight: enabled ? FontWeight.bold : FontWeight.normal,
              color: enabled ? Colors.indigo[800] : Colors.grey[600],
            ),
          ),
          const SizedBox(width: 4),
          Icon(
            enabled ? Icons.check_circle : Icons.cancel,
            size: 14,
            color: enabled ? Colors.green[600] : Colors.grey[400],
          ),
        ],
      ),
    );
  }

  Widget _buildAutoSelectionItem(String label, String value, Color color) {
    return Container(
      padding: const EdgeInsets.all(8),
      decoration: BoxDecoration(
        color: Colors.white,
        borderRadius: BorderRadius.circular(6),
      ),
      child: Column(
        crossAxisAlignment: CrossAxisAlignment.start,
        children: [
          Text(
            label,
            style: TextStyle(
              fontSize: 10,
              color: Colors.grey[600],
            ),
          ),
          const SizedBox(height: 2),
          Text(
            value,
            style: TextStyle(
              fontSize: 12,
              fontWeight: FontWeight.bold,
              color: color,
            ),
          ),
        ],
      ),
    );
  }

  Widget _buildValidationSummaryCard(ModelInfo modelInfo) {
    final validation = modelInfo.validationSummary;
    final coverage = modelInfo.dataCoverage;

    return Container(
      padding: const EdgeInsets.all(12),
      decoration: BoxDecoration(
        color: Colors.indigo[50],
        borderRadius: BorderRadius.circular(12),
        border: Border.all(color: Colors.indigo[100]!),
      ),
      child: Column(
        crossAxisAlignment: CrossAxisAlignment.start,
        children: [
          Row(
            children: [
              Icon(Icons.verified, color: Colors.indigo[700], size: 18),
              const SizedBox(width: 8),
              const Text(
                '验证与数据覆盖',
                style: TextStyle(fontSize: 14, fontWeight: FontWeight.bold),
              ),
            ],
          ),
          const SizedBox(height: 10),
          if (validation != null) ...[
            Row(
              children: [
                Expanded(
                  child: _buildModelStatItem(
                    Icons.rule,
                    '验证方式',
                    validation.method ?? 'N/A',
                    Colors.indigo[800]!,
                  ),
                ),
                const SizedBox(width: 12),
                Expanded(
                  child: _buildModelStatItem(
                    Icons.repeat,
                    '折数',
                    validation.cvFolds?.toString() ?? '—',
                    Colors.deepPurple[700]!,
                  ),
                ),
              ],
            ),
            const SizedBox(height: 10),
            Row(
              children: [
                Expanded(
                  child: _buildModelStatItem(
                    Icons.assessment,
                    'CV MAE',
                    validation.cvMaeMean != null
                        ? '${validation.cvMaeMean!.toStringAsFixed(2)} kW ± ${validation.cvMaeStd?.toStringAsFixed(2) ?? "0"}'
                        : 'N/A',
                    Colors.teal[700]!,
                  ),
                ),
                const SizedBox(width: 12),
                Expanded(
                  child: _buildModelStatItem(
                    Icons.check_circle,
                    'Holdout MAE',
                    validation.holdoutMae != null
                        ? '${validation.holdoutMae!.toStringAsFixed(2)} kW'
                        : 'N/A',
                    Colors.blueGrey[700]!,
                  ),
                ),
              ],
            ),
            const SizedBox(height: 10),
          ],
          if (coverage != null)
            Container(
              padding: const EdgeInsets.all(12),
              decoration: BoxDecoration(
                color: Colors.white,
                borderRadius: BorderRadius.circular(8),
                border: Border.all(color: Colors.indigo[100]!),
              ),
              child: Column(
                crossAxisAlignment: CrossAxisAlignment.start,
                children: [
                  Row(
                    children: [
                      Icon(Icons.date_range, size: 16, color: Colors.indigo[700]),
                      const SizedBox(width: 6),
                      Text(
                        '数据覆盖区间',
                        style: TextStyle(
                          fontSize: 12,
                          fontWeight: FontWeight.bold,
                          color: Colors.indigo[800],
                        ),
                      ),
                    ],
                  ),
                  const SizedBox(height: 6),
                  Text(
                    '${coverage.start ?? "N/A"}  至  ${coverage.end ?? "N/A"}',
                    style: TextStyle(
                      color: Colors.grey[800],
                      fontWeight: FontWeight.w600,
                    ),
                  ),
                  const SizedBox(height: 4),
                  Text(
                    '跨度: ${coverage.spanDays != null ? "${coverage.spanDays} 天" : "未知"} · 样本: ${coverage.rows ?? 0}',
                    style: TextStyle(fontSize: 12, color: Colors.grey[700]),
                  ),
                ],
              ),
            ),
        ],
      ),
    );
  }

  Widget _buildSolverDiagnosticsCard(OptimizationData optimization) {
    final diag = optimization.diagnostics;
    final hits = optimization.constraintHits;
    if (diag == null && hits == null) {
      return const SizedBox.shrink();
    }

    return Card(
      elevation: 3,
      shape: RoundedRectangleBorder(borderRadius: BorderRadius.circular(12)),
      child: Padding(
        padding: const EdgeInsets.all(16.0),
        child: Column(
          crossAxisAlignment: CrossAxisAlignment.start,
          children: [
            Row(
              children: [
                Icon(Icons.speed, color: Colors.blue[700]),
                const SizedBox(width: 8),
                const Text(
                  '求解器健康度',
                  style: TextStyle(fontSize: 16, fontWeight: FontWeight.bold),
                ),
              ],
            ),
            const SizedBox(height: 12),
            Row(
              children: [
                Expanded(
                  child: _buildModelStatItem(
                    Icons.timer,
                    '求解耗时',
                    diag?.runtimeLabel ?? 'N/A',
                    Colors.blue[700]!,
                  ),
                ),
                const SizedBox(width: 12),
                Expanded(
                  child: _buildModelStatItem(
                    Icons.data_usage,
                    'MIP Gap',
                    diag?.mipGap != null ? '${(diag!.mipGap! * 100).toStringAsFixed(2)}%' : 'N/A',
                    Colors.deepOrange[700]!,
                  ),
                ),
              ],
            ),
            const SizedBox(height: 12),
            Row(
              children: [
                Expanded(
                  child: _buildModelStatItem(
                    Icons.account_tree,
                    'Node',
                    diag?.nodeCount?.toString() ?? '—',
                    Colors.teal[700]!,
                  ),
                ),
                const SizedBox(width: 12),
                Expanded(
                  child: _buildModelStatItem(
                    Icons.loop,
                    '迭代',
                    diag?.iterCount?.toString() ?? '—',
                    Colors.indigo[700]!,
                  ),
                ),
              ],
            ),
            if (hits != null) ...[
              const SizedBox(height: 12),
              Row(
                children: [
                  Expanded(
                    child: _buildModelStatItem(
                      Icons.battery_alert,
                      'SOC 下限命中',
                      '${hits.socMinHits} 次',
                      Colors.red[600]!,
                    ),
                  ),
                  const SizedBox(width: 12),
                  Expanded(
                    child: _buildModelStatItem(
                      Icons.battery_full,
                      'SOC 上限命中',
                      '${hits.socMaxHits} 次',
                      Colors.green[700]!,
                    ),
                  ),
                ],
              ),
              const SizedBox(height: 12),
              Row(
                children: [
                  Expanded(
                    child: _buildModelStatItem(
                      Icons.flash_on,
                      '充电功率封顶',
                      '${hits.maxChargeHits} 小时',
                      Colors.orange[700]!,
                    ),
                  ),
                  const SizedBox(width: 12),
                  Expanded(
                    child: _buildModelStatItem(
                      Icons.bolt,
                      '放电功率封顶',
                      '${hits.maxDischargeHits} 小时',
                      Colors.blueGrey[700]!,
                    ),
                  ),
                ],
              ),
            ],
          ],
        ),
      ),
    );
  }

  Widget _buildModelStatItem(IconData icon, String label, String value, Color color) {
    return Container(
      padding: const EdgeInsets.all(12),
      decoration: BoxDecoration(
        color: Colors.white,
        borderRadius: BorderRadius.circular(8),
        border: Border.all(color: Colors.grey[300]!),
      ),
      child: Column(
        crossAxisAlignment: CrossAxisAlignment.start,
        children: [
          Row(
            children: [
              Icon(icon, color: color, size: 16),
              const SizedBox(width: 6),
              Expanded(
                child: Text(
                  label,
                  style: const TextStyle(
                    fontSize: 11,
                    color: Colors.grey,
                  ),
                  overflow: TextOverflow.ellipsis,
                ),
              ),
            ],
          ),
          const SizedBox(height: 6),
          Text(
            value,
            style: TextStyle(
              fontSize: 13,
              fontWeight: FontWeight.bold,
              color: color,
            ),
            maxLines: 2,
            overflow: TextOverflow.ellipsis,
          ),
        ],
      ),
    );
  }

  /// 5. 关键指标卡片 - 增强版（方案一核心展示）
  Widget _buildKeyMetrics(OptimizationData optimization) {
    final summary = optimization.summary;
    
    // 计算与之前结果的差异（如果有）
    double? savingsDiff;
    if (_previousResult?.optimization != null) {
      final prevSavings = _previousResult!.optimization!.summary.savings;
      savingsDiff = summary.savings - prevSavings;
    }
    
    return Column(
      crossAxisAlignment: CrossAxisAlignment.start,
      children: [
        Row(
          children: [
            Icon(Icons.analytics, color: Colors.blue[700], size: 24),
            const SizedBox(width: 8),
            const Text(
              ' 优化效果',
              style: TextStyle(
                fontSize: 20,
                fontWeight: FontWeight.bold,
              ),
            ),
            const Spacer(),
            // 显示与之前结果的对比
            if (savingsDiff != null && savingsDiff.abs() > 0.01) ...[
              Container(
                padding: const EdgeInsets.symmetric(horizontal: 8, vertical: 4),
                decoration: BoxDecoration(
                  color: savingsDiff > 0 ? Colors.green.withOpacity(0.2) : Colors.red.withOpacity(0.2),
                  borderRadius: BorderRadius.circular(12),
                ),
                child: Row(
                  mainAxisSize: MainAxisSize.min,
                  children: [
                    Icon(
                      savingsDiff > 0 ? Icons.trending_up : Icons.trending_down,
                      size: 14,
                      color: savingsDiff > 0 ? Colors.green[700] : Colors.red[700],
                    ),
                    const SizedBox(width: 4),
                    Text(
                      '${savingsDiff > 0 ? "+" : ""}${savingsDiff.toStringAsFixed(2)}元 vs 上次',
                      style: TextStyle(
                        fontSize: 11,
                        fontWeight: FontWeight.bold,
                        color: savingsDiff > 0 ? Colors.green[700] : Colors.red[700],
                      ),
                    ),
                  ],
                ),
              ),
            ],
          ],
        ),
        const SizedBox(height: 12),
        
        // 三个指标卡片 - 响应式布局
        LayoutBuilder(
          builder: (context, constraints) {
            final isMobile = ResponsiveHelper.isMobile(context);
            
            if (isMobile) {
              // 移动端：垂直排列
              return Column(
                children: [
                  _buildMetricCard(
                    icon: Icons.savings,
                    iconColor: Colors.green,
                    label: '节省金额',
                    value: summary.savingsFormatted,
                    backgroundColor: Colors.green[50]!,
                    valueColor: Colors.green[700]!,
                  ),
                  const SizedBox(height: 12),
                  _buildMetricCard(
                    icon: Icons.percent,
                    iconColor: Colors.orange,
                    label: '节省比例',
                    value: summary.savingsPercentFormatted,
                    backgroundColor: Colors.orange[50]!,
                    valueColor: Colors.orange[700]!,
                  ),
                ],
              );
            } else {
              // 平板/桌面端：水平排列
              return Row(
                children: [
                  Expanded(
                    child: _buildMetricCard(
                      icon: Icons.savings,
                      iconColor: Colors.green,
                      label: '节省金额',
                      value: summary.savingsFormatted,
                      backgroundColor: Colors.green[50]!,
                      valueColor: Colors.green[700]!,
                    ),
                  ),
                  const SizedBox(width: 12),
                  Expanded(
                    child: _buildMetricCard(
                      icon: Icons.percent,
                      iconColor: Colors.orange,
                      label: '节省比例',
                      value: summary.savingsPercentFormatted,
                      backgroundColor: Colors.orange[50]!,
                      valueColor: Colors.orange[700]!,
                    ),
                  ),
                ],
              );
            }
          },
        ),
        
        const SizedBox(height: 12),
        
        // 成本对比卡片
        Card(
          elevation: 2,
          shape: RoundedRectangleBorder(borderRadius: BorderRadius.circular(12)),
          child: Padding(
            padding: const EdgeInsets.all(16.0),
            child: Column(
              crossAxisAlignment: CrossAxisAlignment.start,
              children: [
                Row(
                  children: [
                    Icon(Icons.compare_arrows, color: Colors.blue[700]),
                    const SizedBox(width: 8),
                    const Text(
                      '成本对比',
                      style: TextStyle(
                        fontSize: 16,
                        fontWeight: FontWeight.bold,
                      ),
                    ),
                  ],
                ),
                const SizedBox(height: 16),
                
                _buildCostComparisonRow(
                  '无电池成本',
                  summary.totalCostWithoutBattery,
                  Colors.grey,
                ),
                const SizedBox(height: 8),
                _buildCostComparisonRow(
                  '有电池成本',
                  summary.totalCostWithBattery,
                  Colors.blue,
                ),
                
                const Divider(height: 24),
                
                Row(
                  mainAxisAlignment: MainAxisAlignment.spaceBetween,
                  children: [
                    const Text(
                      '总计节省',
                      style: TextStyle(
                        fontSize: 16,
                        fontWeight: FontWeight.bold,
                      ),
                    ),
                    Text(
                      summary.savingsFormatted,
                      style: TextStyle(
                        fontSize: 20,
                        fontWeight: FontWeight.bold,
                        color: Colors.green[700],
                      ),
                    ),
                  ],
                ),
              ],
            ),
          ),
        ),
      ],
    );
  }

  Widget _buildMetricCard({
    required IconData icon,
    required Color iconColor,
    required String label,
    required String value,
    required Color backgroundColor,
    required Color valueColor,
  }) {
    return Card(
      elevation: 2,
      color: backgroundColor,
      shape: RoundedRectangleBorder(borderRadius: BorderRadius.circular(12)),
      child: Padding(
        padding: const EdgeInsets.all(16.0),
        child: Column(
          crossAxisAlignment: CrossAxisAlignment.start,
          children: [
            Icon(icon, color: iconColor, size: 32),
            const SizedBox(height: 12),
            Text(
              label,
              style: TextStyle(
                fontSize: 14,
                color: Colors.grey[700],
              ),
            ),
            const SizedBox(height: 4),
            Text(
              value,
              style: TextStyle(
                fontSize: 24,
                fontWeight: FontWeight.bold,
                color: valueColor,
              ),
            ),
          ],
        ),
      ),
    );
  }

  Widget _buildCostComparisonRow(String label, double cost, Color color) {
    return Row(
      mainAxisAlignment: MainAxisAlignment.spaceBetween,
      children: [
        Row(
          children: [
            Container(
              width: 12,
              height: 12,
              decoration: BoxDecoration(
                color: color,
                shape: BoxShape.circle,
              ),
            ),
            const SizedBox(width: 8),
            Text(
              label,
              style: const TextStyle(fontSize: 14),
            ),
          ],
        ),
        Text(
          '¥${cost.toStringAsFixed(2)}',
          style: TextStyle(
            fontSize: 16,
            fontWeight: FontWeight.w600,
            color: color,
          ),
        ),
      ],
    );
  }

  /// 7. 策略详情
  Widget _buildStrategyDetails(OptimizationData optimization) {
    final strategy = optimization.strategy;
    
    return Card(
      elevation: 4,
      shape: RoundedRectangleBorder(borderRadius: BorderRadius.circular(16)),
      child: Padding(
        padding: const EdgeInsets.all(20.0),
        child: Column(
          crossAxisAlignment: CrossAxisAlignment.start,
          children: [
            Row(
              children: [
                Icon(Icons.schedule, color: Colors.blue[700], size: 24),
                const SizedBox(width: 8),
                const Text(
                  ' 充放电策略',
                  style: TextStyle(
                    fontSize: 18,
                    fontWeight: FontWeight.bold,
                  ),
                ),
              ],
            ),
            const SizedBox(height: 16),
            
            // 充电时段
            _buildStrategyRow(
              icon: Icons.battery_charging_full,
              iconColor: Colors.green,
              label: '充电时段',
              value: strategy.chargingHoursFormatted,
              count: '${strategy.chargingCount} 小时',
              backgroundColor: Colors.green[50]!,
            ),
            
            const SizedBox(height: 12),
            
            // 放电时段
            _buildStrategyRow(
              icon: Icons.flash_on,
              iconColor: Colors.red,
              label: '放电时段',
              value: strategy.dischargingHoursFormatted,
              count: '${strategy.dischargingCount} 小时',
              backgroundColor: Colors.red[50]!,
            ),
            
            const SizedBox(height: 16),
            
            // 统计信息
            Container(
              padding: const EdgeInsets.all(12),
              decoration: BoxDecoration(
                color: Colors.blue[50],
                borderRadius: BorderRadius.circular(8),
              ),
              child: Column(
                children: [
                  _buildStatRow(
                    '总充电量',
                    '${optimization.summary.totalCharged.toStringAsFixed(2)} kWh',
                  ),
                  const Divider(height: 16),
                  _buildStatRow(
                    '总放电量',
                    '${optimization.summary.totalDischarged.toStringAsFixed(2)} kWh',
                  ),
                  const Divider(height: 16),
                  _buildStatRow(
                    '循环效率',
                    '${optimization.summary.cycleEfficiency.toStringAsFixed(1)}%',
                  ),
                ],
              ),
            ),
          ],
        ),
      ),
    );
  }

  Widget _buildStrategyRow({
    required IconData icon,
    required Color iconColor,
    required String label,
    required String value,
    required String count,
    required Color backgroundColor,
  }) {
    return Container(
      padding: const EdgeInsets.all(12),
      decoration: BoxDecoration(
        color: backgroundColor,
        borderRadius: BorderRadius.circular(8),
      ),
      child: Column(
        crossAxisAlignment: CrossAxisAlignment.start,
        children: [
          Row(
            children: [
              Icon(icon, color: iconColor, size: 20),
              const SizedBox(width: 8),
              Text(
                label,
                style: const TextStyle(
                  fontSize: 14,
                  fontWeight: FontWeight.bold,
                ),
              ),
              const Spacer(),
              Container(
                padding: const EdgeInsets.symmetric(horizontal: 8, vertical: 4),
                decoration: BoxDecoration(
                  color: iconColor.withOpacity(0.2),
                  borderRadius: BorderRadius.circular(4),
                ),
                child: Text(
                  count,
                  style: TextStyle(
                    fontSize: 12,
                    fontWeight: FontWeight.bold,
                    color: iconColor,
                  ),
                ),
              ),
            ],
          ),
          const SizedBox(height: 8),
          Text(
            value,
            style: TextStyle(
              fontSize: 13,
              color: Colors.grey[700],
            ),
          ),
        ],
      ),
    );
  }

  Widget _buildStatRow(String label, String value) {
    return Row(
      mainAxisAlignment: MainAxisAlignment.spaceBetween,
      children: [
        Text(
          label,
          style: const TextStyle(fontSize: 14),
        ),
        Text(
          value,
          style: const TextStyle(
            fontSize: 14,
            fontWeight: FontWeight.bold,
          ),
        ),
      ],
    );
  }

  /// 空状态
  Widget _buildEmptyState() {
    return Card(
      elevation: 2,
      shape: RoundedRectangleBorder(borderRadius: BorderRadius.circular(12)),
      child: Padding(
        padding: const EdgeInsets.all(40.0),
        child: Column(
          children: [
            Icon(
              Icons.analytics_outlined,
              size: 80,
              color: Colors.grey[400],
            ),
            const SizedBox(height: 16),
            Text(
              '开始优化',
              style: TextStyle(
                fontSize: 20,
                fontWeight: FontWeight.bold,
                color: Colors.grey[700],
              ),
            ),
            const SizedBox(height: 8),
            Text(
              '调整初始电量参数，点击"开始智能调度"按钮',
              textAlign: TextAlign.center,
              style: TextStyle(
                fontSize: 14,
                color: Colors.grey[600],
              ),
            ),
          ],
        ),
      ),
    );
  }

  /// 信息对话框
  void _showInfoDialog() {
    showDialog(
      context: context,
      builder: (context) => AlertDialog(
        title: const Row(
          children: [
            Icon(Icons.info, color: Colors.blue),
            SizedBox(width: 8),
            Text('关于能源优化'),
          ],
        ),
        content: const SingleChildScrollView(
          child: Column(
            crossAxisAlignment: CrossAxisAlignment.start,
            mainAxisSize: MainAxisSize.min,
            children: [
              Text(
                '功能说明',
                style: TextStyle(fontWeight: FontWeight.bold, fontSize: 16),
              ),
              SizedBox(height: 8),
              Text('• 基于机器学习预测未来24小时能源负载'),
              Text('• 使用混合整数规划优化电池充放电策略'),
              Text('• 考虑峰谷电价，最小化总购电成本'),
              SizedBox(height: 16),
              Text(
                '电价时段',
                style: TextStyle(fontWeight: FontWeight.bold, fontSize: 16),
              ),
              SizedBox(height: 8),
              Text('• 谷时 (00:00-08:00, 22:00-24:00): 0.3 元/kWh'),
              Text('• 平时 (08:00-18:00): 0.6 元/kWh'),
              Text('• 峰时 (18:00-22:00): 1.0 元/kWh'),
              SizedBox(height: 16),
              Text(
                '电池参数',
                style: TextStyle(fontWeight: FontWeight.bold, fontSize: 16),
              ),
              SizedBox(height: 8),
              Text('• 容量: 13.5 kWh (Tesla Powerwall)'),
              Text('• 最大功率: 5.0 kW'),
              Text('• 充放电效率: 95%'),
            ],
          ),
        ),
        actions: [
          TextButton(
            onPressed: () => Navigator.of(context).pop(),
            child: const Text('关闭'),
          ),
        ],
      ),
    );
  }
}
\end{lstlisting}

\chapter{测试报告}

\section{测试环境与策略}

\subsection{测试环境}
本系统的测试环境分为开发环境（本地）和生产环境（Google Cloud）。
\begin{itemize}
    \item \textbf{硬件环境}：MacBook Pro (M3 Pro, 18GB RAM) / Google Cloud Run (2 CPU, 4GB RAM)
    \item \textbf{操作系统}：macOS Sequoia 15.0 / Linux (Container Optimized OS)
    \item \textbf{软件依赖}：Python 3.10, Flutter 3.27, Gurobi 11.0, Firebase SDK
\end{itemize}

\subsection{测试策略}
采用金字塔测试策略，包含以下三个层级：
\begin{enumerate}
    \item \textbf{单元测试 (Unit Testing)}：针对后端核心服务（ML、Optimization）和前端独立组件（Models、Utils）进行测试。
    \item \textbf{集成测试 (Integration Testing)}：测试 API 接口的连通性、数据库交互以及前后端数据流转。
    \item \textbf{系统测试 (System Testing)}：模拟完整用户流程，验证定时任务调度、生产环境部署工作流。
\end{enumerate}

\section{测试用例与执行记录}

\subsection{后端单元与集成测试}
使用 \texttt{pytest} 框架执行后端测试，覆盖率约 85\%。主要测试集如下：

\begin{table}[h]
    \centering
    \caption{后端测试用例执行摘要}
    \begin{tabular}{|l|l|l|l|}
        \hline
        \textbf{测试模块}                      & \textbf{测试内容} & \textbf{关键用例}                          & \textbf{结果}           \\ \hline
        \texttt{test\_auth.py}             & 用户认证          & JWT生成/验证、密码哈希                          & \textcolor{green}{通过} \\ \hline
        \texttt{test\_ml\_service.py}      & 机器学习服务        & 数据预处理、模型训练、预测输出格式                      & \textcolor{green}{通过} \\ \hline
        \texttt{test\_optimization.py}     & Gurobi优化      & 约束条件有效性、目标函数收敛、边界值                     & \textcolor{green}{通过} \\ \hline
        \texttt{test\_api\_integration.py} & API 接口        & GET /config, POST /analyze (Mock Auth) & \textcolor{green}{通过} \\ \hline
        \texttt{test\_mlops.py}            & MLOps 流程      & 数据漂移检测、缺失值处理鲁棒性                        & \textcolor{green}{通过} \\ \hline
    \end{tabular}
\end{table}

\subsection{前端测试}
使用 \texttt{flutter test} 进行测试，主要验证 Widget 构建和数据模型解析。

\begin{table}[h]
    \centering
    \caption{前端测试用例执行摘要}
    \begin{tabular}{|l|l|l|}
        \hline
        \textbf{测试文件}              & \textbf{测试内容}                        & \textbf{结果}           \\ \hline
        \texttt{widget\_test.dart} & 基础 Widget 构建、App 启动冒烟测试              & \textcolor{green}{通过} \\ \hline
        \texttt{model\_test.dart}  & JSON 序列化/反序列化 (User, AnalysisResult) & \textcolor{green}{通过} \\ \hline
    \end{tabular}
\end{table}

\subsection{生产工作流验证脚本}
为确保部署稳定性，开发了自动化验证脚本 \texttt{verify\_production\_workflow.py}，模拟真实生产环境的周期性任务。

\begin{lstlisting}[style=codestyle, language=bash, caption={生产环境验证脚本执行日志}]
$ python verify_production_workflow.py
[KEY] 使用凭证: .../service-account-key.json
============================================================
[TEST] 测试任务 1: 每小时数据抓取 (Data Fetching)
============================================================
1. ExternalDataService 初始化成功
2. 执行 fetch_and_publish()...
[OK] 数据抓取成功！数据已保存到 Firebase Storage。

============================================================
[TEST] 测试任务 2: 每日模型训练 (Model Training)
============================================================
1. EnergyPredictor 初始化成功
2. 执行 train_model(use_firebase_storage=True)...
[OK] 模型训练成功！
   MAE: 45.21
   RMSE: 58.90
   模型文件已更新到 Firebase Storage。

============================================================
[ALL PASS] 全部验证通过！您的生产工作流在本地测试正常。
\end{lstlisting}

\section{测试结论}
经过全面的单元测试、集成测试及生产模拟验证，系统表现稳定。
\begin{enumerate}
    \item \textbf{功能符合性}：所有核心功能（数据抓取、负载预测、优化调度、可视化展示）均按照需求规格说明书实现。
    \item \textbf{鲁棒性}：后端服务具备异常捕获机制，能够处理 Gurobi 求解失败、外部 API 超时等异常情况。
    \item \textbf{安全性}：API 接口均受 JWT 验证保护，敏感配置已从代码中分离。
\end{enumerate}

系统已达到上线标准，可以进行生产部署。